% =====================================================================
%  Zusammenfassung EPP (HS26) – LaTeX-Quelle
%  Kompilieren: pdflatex (2x) oder latexmk -pdf main.tex
%  In VS Code mit der Extension "LaTeX Workshop": main.tex öffnen,
%  speichern (Ctrl+S) baut automatisch, Ctrl+Alt+V zeigt das PDF.
%
%  Format: A4 hoch, 2 Spalten, 8pt.
%  Gliederung: nach Themen, nicht nach Vorlesungswochen.
%  Inhalt: eine Datei pro Thema in sections/
% =====================================================================
\documentclass[8pt]{extarticle}

% ---------------------------------------------------------------------
% Pakete
% ---------------------------------------------------------------------
\usepackage[utf8]{inputenc}
\usepackage[T1]{fontenc}
\IfFileExists{ngerman.ldf}{\usepackage[ngerman]{babel}}{%
  \renewcommand{\contentsname}{Inhaltsverzeichnis}}
\usepackage{helvet}
\renewcommand{\familydefault}{\sfdefault}

\usepackage[a4paper,left=1.0cm,right=1.0cm,top=1.0cm,bottom=1.1cm,
            footskip=0.5cm]{geometry}
\usepackage{amsmath,amssymb}
\usepackage{graphicx}
\usepackage{multicol}
\usepackage{booktabs}
\usepackage{array}
\usepackage{tabularx}
\usepackage{environ}        % damit tabularx in \ztab funktioniert
\usepackage{enumitem}
\usepackage[dvipsnames]{xcolor}
\usepackage{titlesec}
\usepackage[most]{tcolorbox}
\usepackage{listings}
\usepackage{fancyhdr}
\usepackage{lastpage}
\usepackage[colorlinks=true,linkcolor=blue!60,urlcolor=blue!60]{hyperref}

\graphicspath{{figures/}}

% ---------------------------------------------------------------------
% Metadaten (hier anpassen)
% ---------------------------------------------------------------------
\newcommand{\doctitle}{Zusammenfassung - EPP}
\newcommand{\docauthors}{Jan Vogel}
\newcommand{\docdate}{\today}
\newcommand{\docversion}{v2.1}

% ---------------------------------------------------------------------
% Layout: 2 Spalten, Abstände
% Platz knapp? Zuerst hier an den Rändern drehen, dann an der
% Schriftgrösse in \documentclass ganz oben.
% ---------------------------------------------------------------------
\setlength{\columnsep}{10pt}
\setlength{\columnseprule}{0.4pt}
\renewcommand{\columnseprulecolor}{\color{black!30}}
\setlength{\parindent}{0pt}
\setlength{\parskip}{0.5pt}
\setlength{\tabcolsep}{3pt}
\setcounter{secnumdepth}{2}
\setcounter{tocdepth}{1}

\setlist[itemize]{nosep,leftmargin=1.1em,topsep=1pt,label=\textcolor{accent}{\textbullet}}
\setlist[enumerate]{nosep,leftmargin=1.5em,topsep=1pt}
\setlist[description]{nosep,topsep=1pt,leftmargin=0pt,font=\normalfont\bfseries}

\setlength{\jot}{5pt}
\allowdisplaybreaks
\AtBeginDocument{%
  \setlength{\abovedisplayskip}{3pt}\setlength{\belowdisplayskip}{3pt}%
  \setlength{\abovedisplayshortskip}{3pt}\setlength{\belowdisplayshortskip}{3pt}}

% ---------------------------------------------------------------------
% Farbe: eine durchgehende Akzentfarbe fuer das ganze Dokument.
% Die Boxenfarben stehen fuer die ART des Inhalts (Falle, Rezept,
% Beispiel) -- nicht fuer eine Vorlesungswoche.
% ---------------------------------------------------------------------
\definecolor{accent}{HTML}{1F4E79}
\colorlet{accentbar}{accent!14}

% ---------------------------------------------------------------------
% Überschriften
% ---------------------------------------------------------------------
\titleformat{\section}[block]{\normalfont}{}{0pt}{\sectionbar}
\titleformat{name=\section,numberless}[block]{\normalfont}{}{0pt}{\sectionbarnonum}
\titleformat{\subsection}[block]{\normalfont}{}{0pt}{\subsectionbar}
\newcommand{\sectionbar}[1]{\colorbox{accentbar}{\parbox{\dimexpr\linewidth-2\fboxsep\relax}{\large\strut\textcolor{accent}{\textbf{\thesection\ #1}}}}}
\newcommand{\sectionbarnonum}[1]{\colorbox{accentbar}{\parbox{\dimexpr\linewidth-2\fboxsep\relax}{\large\strut\textcolor{accent}{\textbf{#1}}}}}
\newcommand{\subsectionbar}[1]{{\color{accent}\rule{2pt}{8pt}}\hskip3pt{\footnotesize\textbf{\thesubsection\ #1}}}

\titlespacing*{\section}{0pt}{3pt}{1pt}
\titlespacing*{\subsection}{0pt}{3pt}{0.5pt}

% ---------------------------------------------------------------------
% Kopf-/Fusszeile
% ---------------------------------------------------------------------
\pagestyle{fancy}
\fancyhf{}
\renewcommand{\headrulewidth}{0pt}
\renewcommand{\footrulewidth}{0.4pt}
\fancyfoot[L]{\tiny \docauthors\ --- Zusammenfassung --- EPP --- \docversion}
\fancyfoot[R]{\tiny \thepage/ \pageref{LastPage}}

% ---------------------------------------------------------------------
% Eigene Umgebungen und Befehle
% ---------------------------------------------------------------------
\newcommand{\R}[1]{\texttt{#1}}          % Register / Code / Befehl
\newcommand{\T}[1]{\textbf{#1}}          % Schluesselbegriff
\newcommand{\ra}{\,$\rightarrow$\,}
\newcommand{\SW}[1]{\textcolor{accent}{\textbf{#1}}\hskip0.4em}  % Stichwort am Zeilenanfang
\newcommand{\merkzeile}{\par\textcolor{accent}{$\blacktriangleright$}\ } % Merksatz-Zeile

% Kompakte Tabelle ueber die Spaltenbreite. Erste Zeile = Kopfzeile.
%   \begin{ztab}{@{}Xr@{}} Kopf & Kopf \\ \midrule ... \end{ztab}
% (\NewEnviron statt \newenvironment, sonst findet tabularx sein Ende nicht)
\NewEnviron{ztab}[1]{%
  \par\footnotesize
  \begin{tabularx}{\linewidth}{#1}\toprule\BODY\bottomrule\end{tabularx}%
  \par\vspace{1pt}\normalsize}

% Boxen -- die Farbe sagt, WAS fuer ein Inhalt drinsteht
\tcbset{zfbox/.style={enhanced,arc=1.5mm,boxrule=0.6pt,
  left=3pt,right=3pt,top=2pt,bottom=2pt,
  fonttitle=\sffamily\bfseries\footnotesize,fontupper=\footnotesize,
  before skip=2pt,after skip=2pt}}

\newtcolorbox{formel}[1][]{zfbox,colback=accent!4,colframe=accent,
  coltitle=white,colbacktitle=accent,title={#1}}         % Formel/Kernaussage
\newtcolorbox{falle}[1][Falle]{zfbox,colback=red!4,colframe=red!65,
  coltitle=white,colbacktitle=red!65,title={#1}}         % typischer Fehler
\newtcolorbox{rezept}[1][Rezept]{zfbox,colback=green!4,colframe=green!40!black,
  coltitle=white,colbacktitle=green!40!black,title={#1}} % Schritt fuer Schritt
\newtcolorbox{bsp}[1][Beispiel]{zfbox,colback=yellow!8,colframe=black!55,
  coltitle=white,colbacktitle=black!55,title={#1}}       % durchgerechnet
\newtcolorbox{achtung}[1][Achtung]{zfbox,colback=red!4,colframe=red!65,
  coltitle=white,colbacktitle=red!65,title={#1}}

% Bild in Spaltenbreite (Anteil von \linewidth)
\newcommand{\fig}[2][0.9]{\par\vspace{2pt}{\centering\includegraphics[width=#1\linewidth]{#2}\par}\vspace{1pt}}

% Code
\lstset{
  basicstyle=\ttfamily\scriptsize,
  columns=fullflexible, keepspaces=true,
  aboveskip=2pt, belowskip=2pt, xleftmargin=3pt,
  commentstyle=\color{black!55}\itshape,
  morecomment=[l]{\#}, morecomment=[l]{;},
  literate={ä}{{\"a}}1 {ö}{{\"o}}1 {ü}{{\"u}}1 {ß}{{\ss}}1
}

% =====================================================================
\begin{document}

% ------------------------- Kopf + Inhalt -----------------------------
% Kein \newpage nach dem Kopf: eine eigene Titelseite waere zu teuer.
% Alles laeuft in einem Spaltenfluss.
% Reihenfolge aendern: \input-Zeilen umsortieren. Weglassen: auskommentieren.
\begin{multicols*}{2}
  {\centering
   {\Large\bfseries \doctitle\par}\vspace{1pt}
   {\footnotesize \docauthors \quad\textbullet\quad \docdate\par}}
  \vspace{4pt}
  % !TEX root = ../main.tex
% ---------------------------------------------------------------------
%  Navigation: "Wenn die Aufgabe lautet ..." -- von der Aufgabenstellung
%  direkt zur Abschnittsnummer. Steht vor dem Inhaltsverzeichnis.
% ---------------------------------------------------------------------

\begin{tcolorbox}[zfbox, colback=black!2, colframe=black!45,
  coltitle=white, colbacktitle=black!45, title={Wenn die Aufgabe lautet \dots}]
\begin{tabularx}{\linewidth}{@{}Xr@{}}
Adresse in Tag/Index/Offset zerlegen & \ref{rez:zerlegung} \\
Mittlere Zugriffszeit (AMAT) rechnen & \ref{form:amat} \\
Trashing-Szenario konstruieren & \ref{bsp:trash} \\
Page-Table-Grösse 1- vs.\ 2-stufig & \ref{sec:mmu} \\
NEON-Befehl dekodieren & \ref{rez:neondecode} \\
Speedup aus Lanes oder Kernen & \ref{form:amdahl} \\
Warum ist diese Schleife nicht vektorisierbar & \ref{sec:simdgrenzen} \\
PL330-\R{CCR} dekodieren, Bytes zählen & \ref{rez:ccr} \\
\R{TI}-Register / \R{SRC\_INC} setzen & \ref{rez:ti} \\
Bandbreite aus Takt und Busbreite & \ref{form:bw} \\
Warum darf DMA keine virtuelle Adresse & \ref{falle:dma} \\
RISC-V: Addition mit Carry ohne Flags & \ref{sec:rvflags} \\
Race Condition / Mutex / Prio-Inversion & \ref{sec:kohaerenz} \\
CGRA-Latenz und PE-Auslastung & \ref{sec:cgra} \\
Docker-Kommando zu einer Situation & \ref{sec:docker} \\
Image- und Container-Grösse rechnen & \ref{sec:docker} \\
\end{tabularx}
\vspace{2pt}
\textcolor{accent}{\rule{7pt}{7pt}}\,Formel \;
\textcolor{red!65}{\rule{7pt}{7pt}}\,Falle \;
\textcolor{green!40!black}{\rule{7pt}{7pt}}\,Rezept \;
\textcolor{black!55}{\rule{7pt}{7pt}}\,Beispiel \;
-- Abschnitt~11 ist das Stichwortregister.
\end{tcolorbox}

  \vspace{2pt}
  \tableofcontents
  \vspace{3pt}

  % --- Speicher ---
  % !TEX root = ../main.tex
% ---------------------------------------------------------------------
%  Speicherhierarchie und Cache
% ---------------------------------------------------------------------
\section{Speicherhierarchie \& Cache}
\label{sec:cache}

\SW{Grundproblem} Speicher ist schnell \emph{oder} gross, nie beides
(Leitungslänge, Adressdecodierung, Preis/Bit). Jede Stufe:
ca.\ 3--10$\times$ langsamer, 10--100$\times$ grösser.

\begin{ztab}{@{}Xrr@{}}
Ebene (Raspberry Pi 5) & Grösse & Zugriff \\ \midrule
Register & $\sim$100\,B & 1\,Cycle \\
L1 (I und D getrennt) & 64\,kB/Core & 4 \\
L2 & 512\,kB/Core & 11 \\
L3 (shared) & 2\,MB & 35 \\
RAM & 4\,GB & $\sim$100 \\
SD-Karte (Swap) & 128\,GB & $\sim$100\,000 \\
\end{ztab}
Ohne Cache wäre eine 2.4-GHz-CPU faktisch eine 24-MHz-CPU.

\subsection{Locality of Reference}
\begin{description}
\item[Zeitlich] dieselbe Adresse bald wieder (Schleifenvariable, Stack).
\item[Räumlich] Nachbaradresse bald (linearer Code, Arrays).
\end{description}
\SW{Zahlen} 90--95\,\% der Befehle laufen linear weiter. Von den Sprüngen
$\sim$80\,\% kurze if-then-else, $\sim$20\,\% Schleifensprünge (davon
85\,\% genommen). \ra Der Befehlsstrom bleibt in einem kleinen
Adressfenster. Praxis: Hit Rate 66--99\,\%.

\subsection{Cache Line (typisch 64\,B)}
Drei Gründe, nie einzelne Bytes zu cachen:
\begin{enumerate}
\item \T{Latenz amortisieren}: RAM-Zugriff $\sim$100 Cycles, danach Burst fast gratis.
\item \T{Verwaltung}: Tag/Valid/Dirty pro 4\,B wären grösser als die Nutzdaten.
\item nutzt automatisch die räumliche Lokalität.
\end{enumerate}

\subsection{Adresszerlegung}
\label{sec:zerlegung}
\begin{center}\ttfamily\small
|\ \ Tag\ \ |\ Index\ |\ Offset\ |
\end{center}
\begin{description}
\item[Offset] (LSBs) Position \emph{innerhalb} der Line.
\item[Index] welches \emph{Set} -- reines Adressdecodieren, kein Suchen.
\item[Tag] Rest; mitgespeichert, sagt aus welchem RAM-Block die Line stammt.
\end{description}

\begin{rezept}[Rezept: Adresse zerlegen]
\label{rez:zerlegung}
\begin{enumerate}
\item $\text{Offset-Bits}=\log_2(\text{Linegrösse})$
\item $\#\text{Sets}=\dfrac{\text{Cachegrösse}}{\text{Line}\times\text{Assoziativität}}$,\quad
      $\text{Index-Bits}=\log_2(\#\text{Sets})$
\item $\text{Tag-Bits}=\text{Adressbreite}-\text{Index}-\text{Offset}$
\item Probe: $\#\text{Blöcke}\times\#\text{Sets}\times\text{Line}=$ Adressraum
\end{enumerate}
\end{rezept}

\begin{bsp}[Beispiel: L1 des Pi 5]
64\,kB, 64\,B Lines, 4-way.
$\#\text{Sets}=\frac{65536}{64\cdot4}=256$ \ra
6 Offset-Bit, 8 Index-Bit, Rest ($32-14=18$) Tag.
\end{bsp}

\SW{Zugriffsablauf} Index \ra Zeile ansteuern $\;\;$ Tag vergleichen
\emph{und} Valid prüfen $\;\;$ beides ok $=$ \T{Hit}, mit Offset das Wort
herausgreifen $\;\;$ sonst \T{Miss}: Line laden, Tag eintragen, Valid setzen.

\subsection{Mapping-Arten}
\begin{ztab}{@{}lXl@{}}
Art & Prinzip & Kosten \\ \midrule
Direct mapped & Index erzwingt genau \emph{eine} Line & 1 Komparator \\
Voll assoziativ & jeder Block in jede Line, kein Index & alle Tags, CAM \\
$n$-way set-ass. & Index wählt Set, $n$ freie Plätze darin & $n$ Komparatoren \\
\end{ztab}
\SW{CAM} Voll assoziativ braucht Content Addressable Memory: pro Bit
SRAM-Zelle (6\,T) $+$ EXOR-Komparator, Ausgänge auf ein Wired-NOR.
Schnellstmöglicher Vergleich, aber sehr teuer in Fläche und Strom
\ra heute nur noch für Kleinststrukturen (Micro-TLB).

\begin{falle}[Cache Trashing (nur Direct mapped)]
\label{bsp:trash}
Zwei aktive Codestellen mit \emph{gleichem Index}, aber verschiedenem Tag
werfen sich gegenseitig hinaus. Hit Rate $\to0$, der Cache bleibt zu
7/8 leer.\\
\T{Konstruieren:} Abstand $=k\cdot\frac{\text{Cachegrösse}}{\text{Assoziativität}}$.
Direct mapped, 8\,kB, 64\,B Lines \ra Index $=$ A12:A6; die Adressen
\R{0x0000} und \R{0x2000} kollidieren. Abhilfe: Assoziativität erhöhen.
\end{falle}

\subsection{Verwaltungsbits}
\begin{description}
\item[Valid] Nach Reset steht Müll im Cache, der zufällig Tag $0$ haben
  kann. „Unbenutzt“ $\neq$ „Tag $=0$“. Reset setzt Valid auf 0.
\item[Dirty] (nur D-Cache) Cache ist neuer als RAM. Nur dirty Lines müssen
  beim Verdrängen zurückgeschrieben werden.
\end{description}

\subsection{Replacement}
\begin{ztab}{@{}lX@{}}
FIFO & ältester Eintrag; ignoriert aktuelle Nutzung \\
LRU & am längsten nicht \emph{benutzt}; braucht Ordnungsinfo bei jedem Zugriff, ab 16-way teuer \\
Random & gratis, erstaunlich gut, \emph{immun} gegen pathologische Muster \\
NMRU & Random, aber die zuletzt benutzte Line ausgenommen \\
\end{ztab}
LRU und FIFO kann man mit einer ungünstigen Schleife gezielt zum Versagen
bringen, Random nicht.

\subsection{Write-Through vs.\ Write-Back}
\begin{ztab}{@{}lXX@{}}
 & Write-Through & Write-Back \\ \midrule
RAM & immer aktuell & zeitweise veraltet \\
Bandbreite & hoch (jeder Write) & minimal (nur Verdrängen) \\
Multicore/DMA & unproblematisch & braucht Kohärenzprotokoll \\
\end{ztab}
Schleife, die eine Variable 1000$\times$ hochzählt: Write-Through 1000
RAM-Zugriffe, Write-Back genau einer.

\subsection{Harvard}
Getrennte I- und D-Caches, damit Befehlsholen und Datenzugriff pro Cycle
\emph{parallel} laufen (sonst Stall). Der I-Cache braucht kein Dirty-Bit
und keine Write-Logik. Trennung nur auf L1, ab L2 \T{unified}
$=$ modifizierte Harvard-Architektur.

\subsection{Mittlere Zugriffszeit}
\begin{formel}[AMAT -- Average Memory Access Time]
\label{form:amat}
\[\text{AMAT}=t_{\text{Hit}}+\text{MissRate}\times\text{MissPenalty}\]
Mehrstufig von innen nach aussen verschachteln.
\end{formel}

\begin{bsp}[Beispiel: Pi 5, Hitraten 95/80/50\,\%]
\[4+0.05\,\bigl(11+0.20\,(35+0.50\cdot100)\bigr)=5.4\ \text{Cycles}\]
gegenüber 100 Cycles ohne Cache. Die Kurve kippt brutal, sobald L1 von
95 auf 90\,\% fällt.
\end{bsp}

  % !TEX root = ../main.tex
% ---------------------------------------------------------------------
%  Virtual Memory / MMU
% ---------------------------------------------------------------------
\section{Virtueller Speicher \& MMU}
\label{sec:mmu}

\begin{ztab}{@{}lX@{}}
Cortex-M & MCU, $<$200\,MHz, meist kein Cache, Code aus Flash \\
Cortex-R & \T{MPU}: übersetzt \emph{nicht}, prüft nur Rechte (r/w/x) für
  typ.\ 8 Regionen \ra Taskisolation, „no execute“, Zugriffszeit bleibt
  vorhersagbar \\
Cortex-A & \T{MMU}: kann alles der MPU \emph{und} übersetzt Adressen.
  Voraussetzung für Linux. \\
\end{ztab}

\subsection{Drei Probleme, eine Lösung}
\begin{enumerate}
\item zu wenig RAM \item keine Isolation zwischen Prozessen
\item Fragmentierung des freien Speichers
\end{enumerate}
\SW{Idee} Jeder Prozess bekommt die \emph{Illusion} eines eigenen,
zusammenhängenden Adressraums. \T{Pages} (typ.\ 4\,kB) virtuell,
\T{Frames} physisch, \T{Page Table} ordnet zu.
\begin{itemize}
\item Aufeinanderfolgende Pages dürfen in verstreuten Frames liegen
      \ra Fragmentierung erledigt
\item Jeder Prozess hat eigene Tabelle mit disjunkten Frames
      \ra Isolation erledigt
\end{itemize}

\SW{Paging/Swapping} Reicht das RAM nicht, wandert eine selten benutzte
Page auf die Disk, ihr Eintrag wird invalid. Späterer Zugriff
\ra \T{Page Fault} \ra OS holt zurück, verdrängt eine andere, Befehl wird
wiederholt. Dauert $\sim$100\,000 Cycles
\ra \emph{page miss $=$ task switch}.

\subsection{Adressübersetzung}
\SW{Einstufig} 32\,Bit $=$ 20\,Bit Page-Nummer $+$ 12\,Bit Offset
($2^{12}=4$\,kB). Der Offset wird \emph{nicht} übersetzt (Pages und Frames
sind gleich gross). Page-Nummer indiziert die Tabelle, diese liefert
Frame-Nummer $+$ Rechte.

\begin{falle}[Warum einstufig nicht geht]
$2^{20}$ Einträge $\times$ 4\,B $=$ \T{4\,MB Page Table pro Prozess}.
Bei 50 Prozessen 200\,MB nur für Tabellen.
\end{falle}

\SW{Zweistufig} $10+10+12$. Die ersten 10\,Bit indizieren eine
First-Level-Tabelle (1024 Einträge, 4\,kB); jeder Eintrag zeigt auf eine
Second-Level-Tabelle mit wieder 1024 Einträgen. Ein Prozess, der nur
10\,MB benutzt, braucht wenige Second-Level-Tabellen \ra statt 4\,MB nur
ca.\ 20\,kB. Preis: ein zusätzlicher Speicherzugriff pro Übersetzung.

\subsection{TLB (Translation Lookaside Buffer)}
Ohne TLB kostete jeder Speicherzugriff zwei zusätzliche RAM-Zugriffe nur
für die Übersetzung -- dreimal langsamer als ganz ohne MMU. Der TLB ist
der Cache für Übersetzungen: klein, aber hoch assoziativ.

\SW{Warum klein reicht} Ein Eintrag deckt 4\,kB Adressraum ab; 32 Einträge
schon 128\,kB. Mit Locality sitzt ein Programm fast immer in wenigen Pages.

\begin{ztab}{@{}lXl@{}}
Cortex-A72 & Struktur & Einträge \\ \midrule
Micro-TLB D & voll assoziativ & 32 \\
Micro-TLB I & voll assoziativ & 48 \\
L2-TLB & 4-way set-assoziativ & 1024 \\
\end{ztab}

\begin{formel}[Die Analogie, die alles zusammenhält]
\begin{tabularx}{\linewidth}{@{}XX@{}}
\T{Cache} & \T{Virtual Memory} \\
Cache Line & Page / Frame \\
Tag-RAM & Page Table \\
Valid Bit & Valid Bit \\
Offset in der Line & Offset in der Page \\
Miss \ra nachladen & Page Fault \ra einlagern \\
Replacement in HW & Replacement im OS (SW) \\
\end{tabularx}\\[2pt]
Nur die Zeitskalen unterscheiden sich um Faktor $\sim$1000 -- deshalb
entscheidet dort Hardware, hier Software.
\end{formel}

\subsection{\R{vm\_stat} lesen}
\begin{ztab}{@{}lX@{}}
free/active/inactive & RAM-Zustand, Auslagerungskandidaten \\
wired down & nie auslagerbar (Kernel -- wer behandelte sonst den Fault?) \\
translation faults & Page Faults, meist harmlose \emph{soft faults} \\
pageins / pageouts & Disk-Transfers \\
swapins / swapouts & stetig steigend $=$ zu wenig RAM, System thrasht \\
copy-on-write & Fork teilt Frames bis zum ersten Schreibzugriff \\
\end{ztab}
Achtung: Apple Silicon benutzt 16-kB-Pages, nicht 4\,kB wie im Rest der
Vorlesung.

  % --- Datenparallelitaet ---
  % !TEX root = ../main.tex
% ---------------------------------------------------------------------
%  SIMD mit ARM NEON
% ---------------------------------------------------------------------
\section{SIMD mit ARM NEON}
\label{sec:neon}

\SW{Motivation} Colour adjust, 1600$\times$1200 $\approx$ 2\,Mpixel,
3 Additionen pro Pixel:
\begin{ztab}{@{}Xr@{}}
Variante & Instr./Pixel \\ \midrule
naiv: 3 ADD $+$ LD $+$ ST $+$ DEC $+$ BR & $\geq7$ \\
$+$ Load/Store-Multiple, Unrolling 8 (DEC$+$BR geteilt) & $\geq5.25$ \\
$+$ SIMD 8 & $\geq3.25$ \\
\end{ztab}
Rechnung zu 5.25: pro 8 Pixel $24+8+8+1+1=42$ Instruktionen.
Nach SIMD bleiben im Wesentlichen \emph{Load und Store} übrig
\ra \T{Bottleneck RAM}. 6\,MB Bilddaten passen nicht in 2\,MB L3:
$\sim$95\,000 Misses $\times$ 100 Cycles obendrauf.

\subsection{Lanes}
128-Bit-Register, je nach Befehlssuffix geteilt in
$16\times8$ / $8\times16$ / $4\times32$ / $2\times64$ Bit.
\T{Lane 0 $=$ LSBs}, ganz rechts.
Hardware: Carry-Kette an den Lane-Grenzen auftrennen -- fast gratis,
deshalb ist SIMD so attraktiv.

\begin{formel}[Merksatz]
Der Speedup skaliert mit der Lane-Zahl \ra immer den \emph{kleinsten
ausreichenden} Datentyp wählen. 8\,Bit gibt 16 Lanes, 32\,Bit nur 4.
Darum INT8 in Bildverarbeitung und KI.
\end{formel}

\subsection{Registerfile}
\SW{ARMv7} $16\times$\T{Q} (128\,Bit), $32\times$\T{D} (64),
$32\times$\T{S} (32, VFP, teilt sich Q0--Q7 / D0--D15).
Alle auf denselben Flip-Flops -- \emph{wie eine C-Union}.
\SW{AArch64} (Pi 4/5 mit 64-Bit-OS) $32\times$\T{V}, 128\,Bit, kein
Paar-Aliasing, Syntax \R{ADD V0.4S,V1.4S,V2.4S}.

\begin{falle}[Zwei Fallen beim Registerfile]
1.~Schreiben in \R{D3} zerstört die obere Hälfte von \R{Q1}.\\
2.~NEON hat ein \emph{eigenes} Registerfile und eine eigene Pipeline
\ra Transfers ARM\,$\leftrightarrow$\,NEON kosten viele Cycles. Code, der
ständig hin- und herspringt, ist \emph{langsamer} als reiner Skalarcode.
\end{falle}

\subsection{Datentypen}
\begin{ztab}{@{}lXX@{}}
 & 8/16/32/64 & Bemerkung \\ \midrule
U / S & \R{U8}\dots\R{S64} & unsigned / signed \\
I & \R{I8}\dots\R{I64} & vorzeichen\emph{un}abhängig \\
F & \R{F16}, \R{F32} & \T{kein F64} in NEON! \\
P & \R{P8}, \R{P16} & Polynom über GF(2), carry-less \\
\end{ztab}
\SW{Warum „I“} Addition ist bitweise identisch für signed/unsigned
\ra \R{VADD.I16}. Sättigung ist es \emph{nicht}
\ra nur \R{VQADD.S16}/\R{.U16}. Die Existenz von „I“ verrät also, dass der
Befehl vorzeichenunabhängig ist.
\SW{P} für CRC, AES-GCM, Reed-Solomon (XOR statt Addition, kein Übertrag).
\SW{Float} NEON-Floats sind auf ARMv7 nicht voll IEEE-754 (flush-to-zero)
\ra GCC vektorisiert Float-Schleifen nur mit \R{-ffast-math}.

\subsection{Befehlssyntax dekodieren}
\begin{formel}[Aufbau jedes NEON-Befehls]
\label{rez:neondecode}
\begin{center}\ttfamily\small
V\{mod\}op\{shape\}\{cond\}.\{dt\}\ dest,\ src1,\ src2
\end{center}
\begin{tabularx}{\linewidth}{@{}clX@{}}
\multicolumn{3}{@{}l}{\T{mod}}\\
Q & saturating & klemmt statt umzuwickeln ($250{+}10=255$) \\
H & halving & $(a{+}b)/2$ ohne Überlauf \\
D & doubling & Q15-Festkomma: 1 Bit nach links \\
R & rounding & statt abschneiden, vermeidet Bias \\[2pt]
\multicolumn{3}{@{}l}{\T{shape} (Ziel $\neq$ Quellformat)}\\
-- & & D,\,D \ra D \\
L & long & D,\,D \ra Q \\
W & wide & Q,\,D \ra Q \\
N & narrow & Q,\,Q \ra D \\
\end{tabularx}
\end{formel}

\begin{bsp}[Beispiel: \R{VQRDMULH.S16 Q0,Q1,Q2}]
V$|$\,Q\,R\,D\,$|$MULH$|$.S16\\
$=$ Saturating $+$ Rounding $+$ Doubling Multiply, High half,
8 Lanes à 16\,Bit signed. Der Arbeitsbefehl jedes FIR-Filters --
gelesen, ohne ihn zu kennen.
\end{bsp}

\SW{Standardmuster} 8\,Bit laden \ra mit \T{L} auf 16\,Bit erweitern
\ra in 16\,Bit rechnen \ra mit \T{N} und Sättigung zurück auf 8\,Bit.

\subsection{Befehlsklassen}
\begin{description}
\item[Load/Store] \R{VLDn}/\R{VSTn} ($n{=}1..4$) laden
  \emph{de-interleaved}: \R{VLD3.8} liest 24\,B und sortiert R/G/B in drei
  Register. \emph{Ohne das wäre NEON für Bildverarbeitung nutzlos.}
  Dazu \R{VLDM/VSTM}, \R{VLDR/VSTR}, \R{VPUSH/VPOP}.
\item[Data Processing] \R{VDUP} (Skalar in alle Lanes),
  \R{VMOV} (Immediate), \R{VZIP/VUZP}, \R{VTRN} ($4{\times}4$
  transponieren), \R{VTBL/VTBX} (beliebige Permutation),
  \R{VEXT} (gleitendes Fenster über zwei Register),
  \R{VCVT} (auch Fixed-Point), \R{VREV16}.
\item[Shift] Es gibt \emph{keine} SIMD-Division, darum sind Shifts
  zentral: \R{VSHL/VSHR}, \R{VSRA} (shift $+$ accumulate),
  \R{VSHLL}, \R{VQRSHRN} (Sat$+$Round$+$Shift$+$Narrow -- der
  Standard-Abschlussbefehl), \R{VSLI/VSRI} (Bitfelder packen).
\item[Logical/Compare] Prädikation statt Branch:
  \R{VCEQ/VCGE/VCGT} liefern eine \emph{Maske} (alle 1 / alle 0),
  \R{VBSL} (Bit Select) wählt bitweise \ra das \R{?:} der Vektorwelt.
  Preis: beide Zweige werden gerechnet.
\item[Arithmetik] \R{VABD/VABA} (SAD -- Bewegungsschätzung im Videocodec),
  \R{VMAX/VMIN}, \R{VPADD/VPMAX} (\emph{pairwise} $=$ benachbarte Lanes),
  \R{VCLZ/VCNT}, \R{VRECPE/VRECPS} und \R{VRSQRTE/VRSQRTS}
  (Newton-Raphson statt Division/Wurzel).
\item[Multiply] \R{VMLA/VMLS} (MAC -- Grundoperation von FIR,
  Skalarprodukt, Faltung, Matrixmultiplikation), \R{VFMA}
  (\emph{fused}: nur eine Rundung, numerisch genauer),
  „by scalar“-Varianten, \R{VQDMULH} (Q15).
\end{description}

\subsection{Intrinsics}
\begin{ztab}{@{}lX@{}}
\R{v<op>q\_<dt>} & \R{q} $=$ Quadword (128\,Bit), ohne \R{q} $=$ 64 \\
\R{<typ><br>x<anz>\_t} & \R{uint8x16\_t}, \R{float32x4\_t} \\
\R{...x<länge>\_t} & \R{uint8x8x3\_t} $=$ 3 Register (für \R{VLD3}) \\
\end{ztab}
Header \R{<arm\_neon.h>}. Registerzuteilung und Scheduling macht weiterhin
der Compiler.

\SW{Matrixmultiplikation $4\times4$} Daten \T{column-major}
(\R{C[4*j+i]}). Nicht Skalarprodukt Zeile$\times$Spalte (bräuchte
horizontale Summen), sondern jede Spalte von C als \T{Linearkombination
der Spalten von A}: \R{C0 += A0*B0[0] + A1*B0[1] + \dots} über
\R{vfmaq\_laneq\_f32}. \ra Alle 4 Lanes arbeiten durchgehend, \emph{keine}
horizontale Operation, keine Transposition. 12 von 16 Q-Registern belegt --
Registerdruck ist die harte Grenze ($8\times8$ ginge nicht mehr).

\SW{MAC-/FIR-Filter} In der Schleife nur \emph{vertikal} rechnen
(\R{vmla.f32 q0,q1,q2} lane-weise; \R{!} $=$ Post-Increment spart die
Adressrechnung). Am \emph{Ende} einmal horizontal reduzieren:
\R{vpadd} in $\log_2 n$ Schritten
$[a,b],[c,d]\to[a{+}b,\,c{+}d]\to$ Gesamtsumme.

\begin{falle}[Horizontal nie in der Schleife]
Pairwise/horizontale Operationen brechen die Lane-Isolation und sind
langsam. Ein Akkumulator \emph{pro Lane}, Reduktion genau einmal am Schluss.
\end{falle}

\subsection{Grenzen}
\label{sec:simdgrenzen}
\begin{formel}[Amdahl für SIMD]
\label{form:amdahl}
\[S=\frac{1}{(1-p)+p/s}\qquad p=\text{vektorisierbar},\ s=\text{Lanes}\]
$p{=}0.8,\ s{=}4\Rightarrow2.5\times$; $s{=}16\Rightarrow$ nur $4\times$.
Die 20\,\% Skalarcode deckeln alles.
\end{formel}

\begin{description}
\item[Memory bound] \R{c[i]=a[i]+b[i]}: 1 Rechenoperation auf 3
  Speicherzugriffe. \emph{Zuerst} Datenlayout und Cache-Lokalität,
  \emph{dann} SIMD.
\item[Hindernisse] loop-carried dependency (\R{a[i]=a[i-1]+b[i]});
  Pointer-Aliasing \ra \R{restrict}; Kontrollfluss \ra Masken;
  Float-Assoziativität \ra \R{-ffast-math}.
\item[Tail] 1003 Elemente bei 4 Lanes $=$ 250 Vektoriterationen $+$ 3
  skalar. SIMD lohnt erst ab einigen Dutzend Elementen.
\item[Alignment] unit stride ist Pflicht; Gather zerstört die Effizienz,
  weil jede Lane eine eigene Cache Line anfasst.
\end{description}

\subsection{NEON einschalten (bare metal)}
Nach Reset aus -- wegen Strom (clock gating) und lazy Context Switch des OS.
\T{Zwei} Schalter nötig:
\begin{lstlisting}
MRC p15,0,r0,c1,c0,2   ; CPACR lesen
ORR r0,r0,#0x00f00000  ; cp10+cp11 full access
MCR p15,0,r0,c1,c0,2   ; CPACR schreiben
ISB                    ; Pipeline neu holen!
MOV r0,#0x40000000     ; Bit30 = EN
VMSR FPEXC,r0          ; Einheit selbst einschalten
\end{lstlisting}
CPACR ist die Tür des ARM-Kerns zum Coprozessor-Interface,
\R{FPEXC.EN} der Hauptschalter der Einheit. Ohne \R{ISB} ist das
Verhalten undefiniert.
\SW{Flags} \R{armcc -O3 -Otime --vectorize} (veraltet) bzw.
\R{gcc -O3 -ftree-vectorize}; Diagnose \R{-fopt-info-vec-missed}.

  % !TEX root = ../main.tex
% ---------------------------------------------------------------------
%  RISC-V Vektorerweiterung
% ---------------------------------------------------------------------
\section{SIMD mit RISC-V Vektoren}
\label{sec:rvv}

\begin{formel}[Der eine Unterschied zu NEON]
Die Vektorlänge steht \emph{nicht} im Befehlssatz. Derselbe Binärcode
läuft auf 2-, 4- oder 8-Lane-Hardware -- er braucht nur mehr oder weniger
Takte. Keine Neukompilierung, kein Multiversioning.
\end{formel}

\begin{itemize}
\item 32 Vektorregister \R{v0}--\R{v31}; ein Vektor hält Skalar, Vektor
      oder Matrix
\item \T{VL-Register} definiert die benutzte Länge, $\text{VL}\le\text{MAXVL}$
\item \R{vfadd.s v0,v1,v2}: für $i<\text{VL}$ addieren,
      für $\text{VL}\le i<\text{MAXVL}$ Nullen schreiben
\item \R{vconfig} setzt Datenbreite und Registerzahl und schaltet die
      Einheit ab (Strom!)
\end{itemize}

\subsection{Strip Mining: Tail-Problem in Hardware gelöst}
\begin{lstlisting}
vconfig 0x63        # 32b Daten, 4 vregs, VL=MAXVL
vfld.s  v1, 0(x4)   # Konstante laden + broadcast
loop:
  vsetvl x2, x1     # VL = x2 = min(MAXVL, Rest x1)
  vflw   v0, 0(x3)  # laedt genau VL Elemente
  vfadd.s v2, v1, v0
  vsw    v2, 0(x3)
  slli x5,x2,2      # Bytes = VL * sizeof(float)
  add  x3,x3,x5     # Zeiger weiter
  sub  x1,x1,x2     # Rest verringern
  bne  x1,x0,loop
vconfig 0x0         # Vektoreinheit aus
\end{lstlisting}
\R{vsetvl} liefert $\min(\text{MAXVL},\text{Rest})$ \ra die letzte
Iteration verarbeitet automatisch nur den Rest. \T{Kein Tail-Code.}

\subsection{Befehlsklassen}
\begin{ztab}{@{}lX@{}}
Load/Store & \R{vlw/vsw} unit-stride \\
 & \R{vls*/vss*} \T{strided} (jedes 2., 3., \dots{} Element) \\
 & \R{vlx*/vsx*} \T{indexed} $=$ Gather/Scatter \\
Integer & \R{vadd}, \R{vmul}, \T{\R{vdiv}}, \R{vsll}, \R{vand},
  \R{vseq/vslt}, \R{vclip} (fixed point) \\
Float & \R{vfadd/sub/mul/div}, \R{vfsqrt}, \R{vfmax/min},
  \R{vfclass}; Suffix \R{.h/.s/.d} $=$ FP16/32/64 \\
MAC & \R{vfmadd}, \R{vfmsub}, widening \R{vfmaddwdn} \\
Convert & \R{vfcvt.<ziel>.<quelle>} int$\leftrightarrow$float,
  float$\leftrightarrow$float \\
Movement & \R{vins}/\R{vext} (GPR$\leftrightarrow$Lane),
  \R{vmerge} (maskengesteuert), \R{vrgather}
  (\R{vd[i]=vs1[vs2[i]]}), \R{vsplat}, \R{vslideup/dwn} \\
Masken & \R{vmfirst}, \R{vmpopc}, \R{vmiota} \\
\end{ztab}

\subsection{Gegenüberstellung}
\begin{ztab}{@{}lXX@{}}
 & NEON & RISC-V V \\ \midrule
Länge & fix 128\,Bit & VL variabel \\
Tail & Hand-Code & \R{vsetvl} \\
Masken & Behelf (\R{VBSL}) & eigene Maskenregister \\
Gather/Scatter & -- & \R{vlx*}/\R{vsx*} \\
Division & Newton-Raphson & \R{vdiv} \\
FP64 & -- & \R{.d} \\
\end{ztab}
Die RVV-Erweiterung behebt also genau die vier Schwachstellen, die bei
NEON in Abschnitt~\ref{sec:simdgrenzen} als Grenzen stehen.

  % --- Datentransport ---
  % !TEX root = ../main.tex
% ---------------------------------------------------------------------
%  DMA
% ---------------------------------------------------------------------
\section{DMA -- Direct Memory Access}
\label{sec:dma}

\SW{Motivation} Ein Interrupt pro Byte kostet $\sim$100 Cycles
Kontextwechsel -- \emph{für ein Byte}. DMA: einmal konfigurieren, \emph{ein}
Interrupt für N Byte.

\begin{bsp}[4\,kB von einer Peripherie-FIFO]
byteweiser IRQ: $4096\times100=409\,600$ Cycles CPU-Zeit\\
DMA: $\sim$200 Setup $+$ 1 IRQ $\approx$ 300 Cycles \ra Faktor $\sim$1400.
\end{bsp}

\begin{formel}[Worin der Gewinn wirklich liegt]
Nicht in der Transferdauer -- die bestimmt die Peripherie -- sondern
darin, dass die CPU in dieser Zeit rechnen oder \emph{schlafen} kann.
Darum gehört DMA ins Fach „Power \emph{and} Performance“.
\end{formel}

\subsection{DMA als zweiter Busmaster}
\begin{ztab}{@{}lX@{}}
\R{BR}/\R{BG} & Bus Request/Grant: CPU schaltet ihre Treiber hochohmig \\
\R{RS}/\R{DS} & CPU adressiert und beschreibt die DMA-Register \\
\R{REQ}/\R{ACK} & Handshake mit dem I/O-Gerät \\
\R{IRQ} & Block fertig \\
\end{ztab}
\begin{formel}[Doppelrolle]
Beim \emph{Konfigurieren} ist der DMA \T{Slave}, beim \emph{Transfer}
\T{Master}. Genau das unterscheidet ihn von jeder anderen Peripherie.
(ARM DMAC: AXI-64-Master für Daten, APB-32-Slave für Register.)
\end{formel}

\SW{Arbitrierung}
\begin{ztab}{@{}lXl@{}}
Modus & Verhalten & CPU-Latenz \\ \midrule
Burst/Block & Bus bis Blockende halten & schlecht \\
Cycle Stealing & 1 Zyklus, freigeben, wieder & gut \\
Transparent & nur ungenutzte Zyklen & keine \\
\end{ztab}

\SW{Richtungen} mem\,\ra\,mem (voller Bustakt);
mem\,\ra\,Peripherie und Peripherie\,\ra\,mem (\emph{von der Peripherie
getaktet}, über full/empty-Flags oder req/ack);
externes Interface\,\ra\,mem mit Paketlänge $+$ IRQ (Ethernet).

\begin{formel}[Dasselbe Pacing-Konzept, drei Namen]
\T{Peripheral Request} (\R{DMAWFP}, PL330) $=$
\T{DREQ} (BCM2835) $=$ \T{TREQ} (RP2350).
\end{formel}

\subsection{Die zwei Fallstricke}
\label{falle:dma}
\begin{falle}[1 -- Cache-Kohärenz]
Write-Back heisst: das RAM ist veraltet, der externe DMA sieht den Cache
nicht.\\
\T{CPU\,\ra\,Peripherie:} Puffer noch dirty im Cache \ra DMA liest alte
Daten \ra vorher \T{clean}.\\
\T{Peripherie\,\ra\,CPU:} CPU trifft die alte Cache-Zeile \ra vorher
\T{invalidate}.\\
Alternativen: nicht cachebarer Puffer, oder HW-Kohärenz
(Zynq \T{ACP}-Port über die Snoop Control Unit).
\end{falle}

\begin{falle}[2 -- Physische Adressen]
Der DMA hat \emph{keine} MMU, keinen TLB. Ein Programm-Pointer ist
virtuell \ra stille Datenkorruption statt Segfault.\\
Und: ein virtuell zusammenhängender Puffer ist physisch fast nie
zusammenhängend (1\,MB $=$ 256 beliebig verteilte Frames).\\
\T{Daraus folgt die gesamte restliche Komplexität}: Control Blocks,
verkettete Listen und Befehlssätze existieren, damit ein Transfer aus
mehreren Stücken bestehen kann (\T{Scatter-Gather}). Alternative:
physisch zusammenhängend allozieren (CMA, VideoCore-Mailbox).
\end{falle}

\subsection{Zynq: Bandbreite}
\begin{formel}[Bandbreite]
\label{form:bw}
\[\text{Bandbreite}=\text{Takt}\times\text{Busbreite in Byte}\]
DMA läuft auf \R{CPU\_2x}. Bei CPU\_6x $=$ 800\,MHz (1:3) also 266\,MHz,
Central Interconnect 64\,Bit:
$266\,\text{MHz}\times8\,\text{B}=\T{2.13\,GB/s}$.
Slave Interconnect nur 32\,Bit \ra die Hälfte.
\end{formel}
Die Register hängen am APB mit \R{CPU\_1x} ($\sim$133\,MHz, nicht
gepipelinet) \ra \emph{ein} Registerwrite kostet Dutzende CPU-Cycles.
Darum lohnt DMA erst ab grossen Blöcken.
Clock-Domain-Übergänge (ASYNC, FIFO) kosten Latenz \ra \T{Burst} --
dasselbe Amortisationsargument wie die Cache Line in P1.

\subsection{ARM PL330: DMA mit eigenem Befehlssatz}
\begin{itemize}
\item AMBA-PrimeCell-IP von ARM; die Engine \emph{holt ihr Programm aus
      dem Speicher} über das AXI-Master-Interface
\item 8 Kanäle $=$ Threads, Round-Robin, gleiche Priorität
\item \T{eigener I-Cache}: 8 Lines à 4 Worte; während eines Line-Loads
      darf ein anderer Kanal aktiv sein
\item MFIFO $128\times64$\,Bit, von allen Kanälen geteilt --
      \emph{entkoppelt} Burstlänge und Breite von Quelle und Ziel
\item APB-Slave non-secure \R{0xF800\_4000} / secure \R{0xF800\_3000}
\end{itemize}

\begin{lstlisting}
DMAMOV CCR, SB4 SS64 DB4 DS64  # Src+Dst: 4 Beats a 64b
DMAMOV SAR, 0x1000             # Quelladresse
DMAMOV DAR, 0x4000             # Zieladresse
DMALP 16
  DMALD                        # 4 x 64b = 32 B laden
  DMAST                        # 32 B speichern
DMALPEND
DMAEND
\end{lstlisting}

\begin{rezept}[Rezept: \R{CCR} dekodieren]
\label{rez:ccr}
\R{S}/\R{D} $=$ Source/Destination, \R{B}$n$ $=$ Burst Length (Beats),
\R{S}$m$ $=$ Burst Size (Bit pro Beat).
\begin{enumerate}
\item Bytes pro \R{DMALD} $=n\times m/8$
\item Gesamtmenge $=$ Bytes pro Burst $\times$ Schleifenzahl
\item Prüfen, ob Quell- und Zielseite gleich viel bewegen
\end{enumerate}
Oben: $4\times64/8=32$\,B, $16\times32=\T{512\,B}$.
\end{rezept}

\begin{falle}[„Loop fits in one cache line (4 I-words)“]
Der Schleifenkörper hat genau 4 Instruktionen \ra einmal geholt,
16$\times$ ausgeführt. Ohne Source-Burst (\R{SB1}) braucht man
$4\times$\R{DMALD} $+$ \R{DMAST} $=$ 6 Instruktionen \ra passt nicht mehr,
die Engine fetcht in \emph{jedem} Durchlauf nach. Zweiter Verlust:
4 AXI-Adressphasen statt einer.\\
\T{Das ist Cache Trashing aus P1, eine Ebene tiefer.}
\end{falle}

\SW{Peripherie-Transfer} \R{DMAWFP} blockiert, bis die Peripherie ihren
Request zieht -- die anderen 7 Kanäle laufen weiter. Der Request trägt
einen \emph{Typ} (\R{DxTYPE}: 00 single, 01 burst); danach laufen Befehle
mit Suffix \R{B} nur bei Burst, mit \R{S} nur bei Single.
\T{Prädikation in der Zeit}, analog zu NEONs Masken im Raum.
\R{DMALPFE} $=$ Loop Forever, Abbruch über \R{DRLAST} der Hardware.
\R{DMASTP} $=$ Store \emph{and notify} (das ACK) \ra die letzte
Transaktion wird aus der Schleife gezogen (15$+$1, Loop Peeling).

\begin{ztab}{@{}lX@{}}
Daten & \R{DMALD/ST}, \R{DMALDP/STP}, \R{DMASTZ} (memset in HW) \\
Fluss & \R{DMALP/LPEND/LPFE}, \R{DMAGO/KILL/END/NOP} \\
Sync & \R{DMAWFP}; \R{DMAWFE}/\R{DMASEV} \emph{Kanal-zu-Kanal}
  \ra CPU-freie Pipelines; \R{DMAFLUSHP} \\
Barrieren & \R{DMARMB} / \R{DMAWMB} \\
Adressen & \R{DMAADDH} / \R{DMAADNH} \ra 2D und Stride \\
\end{ztab}

\begin{formel}[VALID/READY -- das universelle AMBA-Handshake]
Sender zieht \R{VALID}, Empfänger \R{READY}; die Übertragung findet im
Takt statt, in dem \emph{beide} high sind. Keine Seite darf auf die andere
warten, bevor sie selbst setzt (sonst Deadlock). Funktioniert bei
beliebigem Geschwindigkeitsverhältnis.
\end{formel}

\subsection{BCM2835 (Raspberry Pi): Control Blocks}
16 Kanäle, keine ISA, stattdessen eine \T{verkettete Liste} von CBs --
8 Worte à 32\,Bit $=$ 32\,B, 32-B-aligned, als \emph{ein} Burst geholt:
\begin{ztab}{@{}clcl@{}}
\# & Feld & \# & Feld \\ \midrule
0 & \R{TI} & 3 & \R{TXFR\_LEN} \\
1 & \R{SOURCE\_AD} & 4 & \R{STRIDE} (2D) \\
2 & \R{DEST\_AD} & 5 & \R{NEXTCONBK} \\
\multicolumn{4}{@{}l}{6--7 reserved, auf 0 setzen} \\
\end{ztab}
\SW{Start} CB-Adresse in \R{CONBLK\_AD} schreiben.
\SW{Stop} Transfer fertig \emph{und} \R{NEXTCONBK}$=0$ \ra ACTIVE-Bit wird
gelöscht.
\SW{Warum 8 Worte} 32\,B ist eine natürliche Burstgrösse und eine
Zweierpotenz \ra CBs sind array-indizierbar und ragen nie über eine Grenze.
\SW{\R{NEXTCONBK}} erlaubt Scatter-Gather und Ringpuffer (letzter CB zeigt
auf den ersten) komplett ohne CPU.

\begin{rezept}[Rezept: \R{TI}-Register setzen]
\label{rez:ti}
Zeigt die Adresse auf \emph{Speicher}, muss sie hochzählen; zeigt sie auf
ein \emph{Peripherieregister} (FIFO), darf sie sich \textbf{nicht} ändern.
\begin{tabularx}{\linewidth}{@{}Xccc@{}}
Transfer & \R{SRC\_INC} & \R{DEST\_INC} & \R{DREQ} \\
mem \ra mem & 1 & 1 & -- \\
mem \ra Periph. & 1 & \textbf{0} & DEST \\
Periph. \ra mem & \textbf{0} & 1 & SRC \\
\end{tabularx}\\[2pt]
Vergessen $=$ der DMA überschreibt den halben Peripherie-Adressraum.
\end{rezept}

\begin{ztab}{@{}lX@{}}
\R{PERMAP} 20:16 & welches Ready-Signal (1--31) \\
\R{WAITS} 25:21 & Wartezyklen einfügen $=$ Bandbreite drosseln \\
\R{BURST\_LENGTH} 15:12 & Burstlänge \\
\R{SRC/DEST\_WIDTH} & 32 oder 128\,Bit \\
\R{WAIT\_RESP} 3 & auf AXI Write Response warten \\
\R{TDMODE} 1 & 2D-Modus, nutzt \R{STRIDE} \ra Rechteck aus Framebuffer \\
\R{INTEN} 0 & Interrupt freigeben \\
\R{SRC\_IGNORE} 11 & nicht lesen \ra schnelles Nullen (memset) \\
\R{DEST\_IGNORE} 7 & nicht schreiben \ra \emph{Prefetch in den L2} \\
\end{ztab}

\begin{formel}[Panic-Signal $=$ QoS in Hardware]
Droht einer Peripherie-FIFO Under- oder Overflow, erhöht der DMA
\emph{dynamisch} sein AXI-Prioritätslevel. Normalbetrieb niedrige
Priorität (CPU nicht stören), Notfall Vorfahrt. Damit muss man sich nicht
zwischen Durchsatz und CPU-Latenz entscheiden.
\end{formel}

\subsection{RP2350: vier Register pro Kanal}
16 Kanäle, kein CB, keine ISA: \R{READ\_ADDR}, \R{WRITE\_ADDR},
\R{TRANS\_COUNT}, \R{CTRL} -- alle vier zählen während der Ausführung mit
(Fortschritt jederzeit lesbar).
\SW{Dual bus ports} je ein Lese- \emph{und} ein Schreibzugriff pro Takt
\ra 1 Transfer/Takt im Dauerbetrieb.

\begin{formel}[Der RP2350-Trick]
Die DMA-Konfigurationsregister sind normale memory-mapped Register --
also kann \emph{ein DMA-Kanal einen anderen programmieren}. Damit baut man
Scatter-Gather aus einem Design ohne jede Listenunterstützung nach.
\end{formel}

\begin{ztab}{@{}lX@{}}
\R{EN}, \R{HIGH\_PRIORITY} & aktiv, mehr Bandbreite \\
\R{DATA\_SIZE} 3:2 & 8/16/32\,Bit \\
\R{INCR\_READ/WRITE} ($+$\R{\_REV}) & hochzählen / rückwärts \\
\R{RING\_SEL}, \R{RING\_SIZE} & „freeze all but $n$ lower bits“
  $=$ \T{Ringpuffer in HW}, Puffer auf $2^n$ ausgerichtet \\
\R{CHAIN\_TO} 16:13 & anderen Kanal auslösen \\
\R{TREQ\_SEL} 22:17 & DREQ-Nummer oder Timer \\
\R{BSWAP} 24 & Endianness gratis \\
\R{SNIFF\_EN} 25 & \T{CRC während des Transfers} \\
\R{BUSY} 26, Fehlerflags 31:29 & Status \\
\end{ztab}
\SW{Weiter} \R{TRANS\_COUNT.MODE} (\R{0x0} stoppen, \R{0x1} neu laden
$+$ IRQ, \R{0xF} endlos); \R{INTR/INTE/INTF/INTS};
\R{TIMER0-3} (Taktteiler \ra Transfer mit exakt fester Rate ganz ohne
Peripherie, z.\,B.\ 44.1\,kHz an einen DAC);
\R{SNIFF\_CTRL} (CRC-32 IEEE\,802.3 / CRC-16-CCITT / 32-Bit-Summe);
\R{MULTI\_CHAN\_TRIGGER}; \R{CHAN\_ABORT}.

\subsection{Drei Philosophien im Vergleich}
\begin{ztab}{@{}lXXX@{}}
 & PL330 & BCM2835 & RP2350 \\ \midrule
Beschreibung & Programm & CB-Liste & 4 Register \\
Kanäle & 8 & 16 & 16 \\
Verkettung & Programm & \R{NEXTCONBK} & \R{CHAIN\_TO} \\
Schleifen & ja & nein & nein \\
Komplexität & hoch & mittel & niedrig \\
\end{ztab}
Drei Punkte auf dem Spektrum Zustandsmaschine\,$\leftrightarrow$\,Prozessor.
\emph{Alle drei lösen dasselbe Problem} -- wer Falle 2 verstanden hat,
kann sich jede der drei Lösungen herleiten.

  % --- Prozessor ---
  % !TEX root = ../main.tex
% ---------------------------------------------------------------------
%  RISC-V ISA
% ---------------------------------------------------------------------
\section{RISC-V Befehlssatz}
\label{sec:riscv}

\SW{RISC-Philosophie} (Patterson, 1980): alle Befehle gleich gross; ein
Befehl pro Takt; Load/Store-Architektur; Harvard; viele
Mehrzweckregister; kleiner einfacher Kern \ra Platz für grosse Caches.

\SW{RISC-V} Offene, lizenz- und gebührenfreie ISA. UC Berkeley, erste
öffentliche Spezifikation 2011, Foundation seit 2019 in der Schweiz
(Sorge vor Exportbeschränkungen).
\begin{itemize}
\item \T{modular}: Basis \R{RV32I} mit $<$50 Befehlen
      (x86 startete 1978 mit 80 und hatte 2015 1338)
\item Erweiterungen \R{C} compressed, \R{M} mult/div, \R{A} atomics,
      \R{F/D/Q} float, \R{V} Vektor
\item Basis und Standard-Erweiterungen sind \emph{frozen}; viel freier
      Opcode-Raum für eigene Erweiterungen
\end{itemize}

\subsection{Registerfile}
31 Mehrzweckregister à 32\,Bit $+$ \R{x0} \T{hardwired 0}.
\R{x1 ra}, \R{x2 sp}, \R{x3 gp}, \R{x4 tp}, \R{t0--t6} (temporär),
\R{s0--s11} (saved), \R{a0--a7} (Argumente/Rückgabe).
\begin{falle}[Konvention, nicht Hardware]
Die Registerfunktionen sind \emph{reine Konvention} -- es gibt z.\,B.\
keine 32-Bit-Maschinenbefehle für Stackoperationen. Der \R{PC} liegt
\emph{ausserhalb} des Registerfiles.
\end{falle}
\SW{Gegenüber ARM} 13 Mehrzweckregister, \R{PC}/\R{SP}/\R{LR} \emph{im}
Registerfile, Barrel Shifter vor der ALU.

\subsection{RV32I: 47 Befehle}
\begin{itemize}
\item feste Breite 32\,Bit (Bit[1:0] $=$ „11“) $+$ 2 illegale
      (alles 0, alles 1)
\item drei 5-Bit-Registeradressen an \emph{festen} Positionen
      \ra Registerlesen läuft parallel zum Decodieren
\item Immediates sind \emph{immer} vorzeichenerweitert
\end{itemize}
\begin{ztab}{@{}lX@{}}
R & Register-Register \\
I / U & short-lower / long-upper Immediate; I$+$U ergeben eine
  32-Bit-Konstante in jedem Register, auch im PC \\
S & Store, Subtyp \T{B} $=$ bedingte Branches \\
J & unbedingte Jumps \\
\end{ztab}

\subsection{Keine ALU-Flags}
\label{sec:rvflags}
\begin{falle}[Es gibt kein C, O, N, Z]
Statt eines Carry-Flags vergleicht \R{SLTU} (set less than, unsigned)
zwei Register. Eine 64-Bit-Addition braucht darum Zusatzbefehle:
\begin{lstlisting}
add  a0,a2,a4   # Summe low
sltu a2,a0,a2   # Carry: (a0<a2) ? 1 : 0
add  a5,a3,a5   # Summe high
add  a1,a2,a5   # Carry addieren
\end{lstlisting}
\end{falle}
\R{CSR} (Control/Status) sind 64-Bit-Zähler für Systemticks, ausgeführte
Befehle und produktive Taktzyklen.

\subsection{RV32C -- Compressed}
34 Befehle à 16\,Bit, jeder bildet auf \emph{genau einen} 32-Bit-Befehl ab
\ra Programmierer und Compiler müssen nichts davon wissen;
32-Bit-Befehle dürfen 16-Bit-aligned liegen.
\begin{itemize}
\item 50--60\,\% der Befehle sind komprimierbar
      \ra \T{25--30\,\% kleinerer Code}
\item Der 16-Bit-Decoder kostet nur $\sim$400--700 Gatter
      ($\sim$2\,\% eines kleinen Kerns)
\end{itemize}
\begin{bsp}[Benchmark insertion sort]
RV32I 76\,B $\;\;$ RV32I$+$C 52\,B (68\,\%) $\;\;$ ARM-Thumb-2 46\,B (60\,\%).
Die gute Dichte von RV32C kommt vor allem vom \emph{grossen Registerfile}
-- es muss kaum kopiert werden.
\end{bsp}

\subsection{ARM vs.\ RISC-V}
\begin{ztab}{@{}XX@{}}
ARM $+$ & RISC-V $+$ \\ \midrule
Industriestandard & offene ISA \\
dichter Maschinencode & einfaches HW-Design \\
grosse Tool- und IP-Basis & kleine Fläche \\
Barrel Shifter vor ALU & grosses Registerfile (31) \\
PC im Registerfile & erweiterbar \\
viele Adressmodi & keine Royalties \\
viele Sleep-Modi & \\ \midrule
ARM $-$ & RISC-V $-$ \\ \midrule
teuer (Royalties) & HW schwer zu kompilieren \\
grosse Fläche, komplex & schlechtere Codedichte \\
patentgeschützt & keine Flags \ra mehr Befehle \\
\end{ztab}

\subsection{RP2350 / Hazard3}
\R{RV32IMAC} $+$ \R{Zba} (Adressgenerierung),
\R{Zbb/Zbs/Zbkb} (Bitmanipulation, Krypto),
\R{Zcb/Zcmp} (weitere compressed, push/pop), \R{Zicsr}.
Debug-/Machine-/User-Modus; PMP mit 8 Regionen, 32-B-Granularität;
\R{Xh3power} (Power-Management), \R{Xh3irq} (verschachtelte,
priorisierte IRQs).
RV32C: alle 16-Bit-Befehle 1 Takt (ausser Traps, IRQs, push/pop).
Ein Sprung auf einen nicht 32-bit-ausgerichteten 32-Bit-Befehl kostet
einen Takt extra (zwei Buszyklen zum Holen).

  % !TEX root = ../main.tex
% ---------------------------------------------------------------------
%  Multicore
% ---------------------------------------------------------------------
\section{Multicore \& Synchronisation}
\label{sec:multicore}

\begin{formel}[Amdahl für Kerne]
\[S(N)=\frac{1}{(1-p)+\dfrac{p}{N}}\]
$p=80\,\%$ parallel, $N=2$:
$S=\frac{1}{0.2+0.4}=1.67$. Dieselbe Formel wie für SIMD-Lanes
(Abschnitt~\ref{form:amdahl}) -- nur heisst $s$ jetzt $N$.
\end{formel}

\SW{Gründe für mehrere Kerne}
\begin{itemize}
\item Komplexität von HW nach SW verschieben: mehrere kleine schnelle
      RISC-Kerne statt eines grossen superskalaren Out-of-Order-Kerns
      -- aber \emph{zuerst} die Caches vergrössern, falls das Optimum
      noch nicht erreicht ist
\item Lizenzkosten steigen typisch \emph{nicht} linear mit der Kernzahl
\item little-big-Architekturen sparen Strom (z.\,B.\ 2 Effizienz- $+$
      2 Performance-Kerne)
\end{itemize}
\SW{HW-Grenze} RAM-Bandbreite. \SW{Nachteil} Datenkohärenz.

\subsection{Multitasking-OS (FreeRTOS)}
Pro Prozessor läuft genau \emph{ein} Task. Bei jedem System-Tick (z.\,B.\
20\,ms Timer-IRQ) kann der Scheduler nach Priorität wechseln -- technisch
nur durch Austausch des \T{Stackpointers}; jeder Task hat seinen eigenen
Stack. Neue Tasks aus der \emph{ready}-Liste; wartende wandern in die
\emph{blocked}-Liste.
\SW{Vorteil} keine Cycles für Polling, einfache Programmierung.
\SW{Problem} Datenkohärenz und geteilte Ressourcen.

\begin{falle}[Reentranz ist Pflicht]
Nicht reentrant: \R{temp} als \emph{globale} Variable in \R{swap()} --
zwei gleichzeitig aktive Instanzen zerstören sich gegenseitig.
Reentrant: \R{temp} \emph{lokal} (auf dem eigenen Stack) -- dann dürfen
Tasks und Kerne denselben Code teilen.
\end{falle}

\subsection{Sechs Ursachen für Kohärenzprobleme}
\label{sec:kohaerenz}
\begin{ztab}{@{}lX@{}}
Stale Data & Core 1 liest den alten Wert aus seinem Cache \\
Write-Write & beide schreiben gleichzeitig, Ergebnis unbestimmt \\
Read-Modify-Write & beide lesen 10, schreiben 11 -- ein Inkrement weg \\
Non-Atomic Multiword & 64\,Bit als zwei 32-Bit-Writes, Leser sieht Mix \\
Out-of-Order & \R{flag=1} wird vor \R{data=42} sichtbar \\
DMA vs.\ CPU & DMA liest, während die CPU noch schreibt \\
\end{ztab}

\begin{bsp}[Race Condition: \R{counter++}]
\begin{lstlisting}
counter++  ->  LOAD counter / ADD 1 / STORE counter

Core0: LOAD 42      Core1: LOAD 42
       ADD 1               ADD 1
       STORE 43            STORE 43   -> 43 statt 44
\end{lstlisting}
Je $10^6$ Inkremente ergeben \emph{nicht} $2\cdot10^6$.
Lösung: \R{lock(); counter++; unlock();} bzw.\ konzeptionell
\R{atomic\_fetch\_add(\&counter,1)}.
\end{bsp}

\subsection{Kooperation paralleler Prozesse}
\SW{Mutual Exclusion brutal} Scheduler, andere CPUs und Interrupts
abschalten und Caches flushen. \SW{Sauber} Semaphore/Mutex.
\begin{ztab}{@{}lX@{}}
Kommunikation & \T{Mailbox} (blockiert) $|$ \T{Message Queue} \\
Synchronisation & \T{Semaphore} (zählend, z.\,B.\ 3 gleiche Ressourcen)
  $|$ \T{Mutex} $|$ \T{Flags} \\
\end{ztab}

\begin{falle}[Priority Inversion]
Ein niederpriorer Task hält das Semaphor und blockiert damit einen
hochprioren; ein mittelpriorer verdrängt den niederprioren dauernd
\ra der hochpriore wartet indirekt auf einen Task, der nie drankommt.\\
\T{Lösung:} \T{Mutex} mit \emph{Priority Inheritance} -- der Halter erbt
vorübergehend die hohe Priorität.\\
\emph{Mars Pathfinder 1997}: genau dieser Fehler, sporadische Resets;
Fix war Priority Inheritance in VxWorks.
\end{falle}

\SW{SMP} identische CPUs, \emph{ein} OS.
\SW{AMP} eine OS-Instanz pro Kern, braucht HW-Unterstützung für
Semaphore/Mutex.

\subsection{FreeRTOS-Kurzreferenz}
\begin{ztab}{@{}lX@{}}
\R{xTaskCreate(fn,name,stack,par,prio,\&h)} & Task anlegen \\
\R{vTaskStartScheduler()} & Scheduler starten \\
\R{vTaskDelay(ms)} & blockieren (nicht pollen!) \\
\R{xQueueCreate(len,size)} & Queue anlegen \\
\R{xQueueSend/Receive(q,\&v,wait)} & Rückgabe \R{pdPASS} prüfen \\
\R{xMessageBufferSend/Receive} & variable Nachrichtenlänge \\
\R{xSemaphoreCreateMutex()} & Mutex (mit Prio-Vererbung) \\
\R{xSemaphoreTake(m,portMAX\_DELAY)} & kritischen Abschnitt betreten \\
\R{xSemaphoreGive(m)} & verlassen \\
\R{xTaskNotifyWait} / \R{\dots FromISR} & leichtgewichtiges Signal \\
\R{vTaskCoreAffinitySet(h,1<<n)} & Task an Kern $n$ binden \\
\end{ztab}
SMP-Konfiguration: \R{configNUM\_CORES 2},
\R{configRUN\_MULTIPLE\_PRIO\dots}, \R{configUSE\_CORE\_AFFINITY},
\R{configUSE\_PREEMPTION}.

\subsection{RP2350 Multiprozessor-Subsystem}
$2\times$Armv8-M (DSP, FPU, Security) \emph{oder} $2\times$\R{RV32IMAC};
Mischbetrieb über \R{ARCHSEL} möglich, braucht aber zwei Programmimages.
\begin{itemize}
\item symmetrische Bussicht: beide Kerne erreichen alle Ressourcen gleich
      schnell
\item pro Kern zwei 32-Bit-Links zum Systembus (Instruktion $+$
      Load/Store $=$ Harvard)
\item ungenutzte Kerne in Reset mit abgeschaltetem Takt
      \ra \emph{null} dynamische Leistung (CMOS)
\item alle Interrupts gehen an beide Kerne, der interne Controller abonniert
\end{itemize}

\SW{SIO (Single-cycle IO)}
\begin{ztab}{@{}lX@{}}
Spinlocks & 32 Stück, Ein-Bit-Flag je Adresse. \emph{Lesen} $=$ Anspruch
  (Rückgabe $\neq0$ \ra erhalten), \emph{Schreiben} (beliebiger Wert) $=$
  freigeben. Bei Gleichstand gewinnt Core\,0. Nur für sehr kurze
  kritische Abschnitte (Spinning ist teuer). \\
FIFOs & Inter-Processor Mailboxes, 32\,Bit breit, 4 Einträge tief;
  Flags \R{VLD}/\R{RDY}/\R{ROE}/\R{WOF} (letzte zwei sticky); IRQ\,25 \\
Doorbells & 8 Flags je Richtung, ein IRQ pro Kern (26), für
  \emph{kumulative, ungeordnete} Ereignisse \\
Rest & \R{CPUID} (0/1), Integer-Divider, Interpolator,
  RISC-V-Platform-Timer, GPIO fürs Bitbanging; alles doppelt für
  secure/non-secure \\
\end{ztab}

\begin{formel}[Global Exclusive Monitor (empfohlen statt Spinlocks)]
Processing Elements: Core\,0, Core\,1 und \emph{der DMA}.
Reservation Page $=$ beliebige 16-B-ausgerichtete SRAM-Region.
Ein exclusive Write gelingt nur, wenn
(a)~ein exclusive Read desselben PE voranging,
(b)~nur eine Reservation offen ist,
(c)~Adresse, Grösse und Rechte übereinstimmen und
(d)~kein anderes PE seither in dieselbe Granule geschrieben hat.
Sonst wird der Write verworfen und meldet das zurück.
\end{formel}
\begin{lstlisting}
arm_acquire:               lr.w rd,(rs1)
  LDREX   r1,[r0]            ; Reservation setzen
  CMP     r1,#0            sc.w rd,rs2,(rs1)
  STREXEQ r1,r2,[r0]         ; rd=0 ok, !=0 fehlgeschlagen
  CMPEQ   r1,#0
  BNE     acquire
\end{lstlisting}
RISC-V zusätzlich \R{amoadd/amoswap/amoand/\dots} (Extension \R{A}).

  % --- Hardware statt Software ---
  % !TEX root = ../main.tex
% ---------------------------------------------------------------------
%  Rekonfigurierbare Hardware
% ---------------------------------------------------------------------
\section{Rekonfigurierbare Hardware}
\label{sec:reconfig}

\SW{Spektrum} fein\,\ra\,grob, schnell\,\ra\,flexibel:
\begin{ztab}{@{}lXl@{}}
Ebene & Baustein & Eigenschaft \\ \midrule
Gatter & \T{FPGA} & sehr feingranular \\
Sequencer & progr.\ FSM & kaum progr., sehr schnell \\
 & Controller & winziger Befehlssatz (Load, Store, Add,
   Jump, Compare), flexibel, langsam \\
 & \T{PIO} & flexibel \emph{und} schnell \\
 & kleine CPU $+$ RAM & sehr flexibel, sehr langsam \\
Operation & \T{CGRA} & routbare ALU/FPU-Blöcke, Datenpfad, VLIW \\
\end{ztab}

\subsection{FPGA}
\begin{itemize}
\item \T{LUT} (Look Up Table) realisiert \emph{jede} Logikfunktion ihrer
      Eingänge; im Beispiel 3, real 4--8
\item Konfigurationsspeicher $=$ SRAM-Zellen in Schieberegister-Anordnung
      (Wiederverwendung der Scan-Test-Struktur)
\item Switch-Matrix-Routing mit Sechs-Transistor-Kreuzungen
\item konfiguriert in VHDL
\end{itemize}
\begin{formel}[Die Zahl, die man sich merkt]
\T{Routing-Delay $>$70\,\% des Gesamt-Delays.} Nicht die Logik ist
langsam, sondern die Verdrahtung.
\end{formel}

\subsection{CGRA}
Array aus \T{PE} (Processing Elements) mit \T{FU} (spezialisierte
ALU/FPU: AND, SHL, ADD, MUL, MAC), typisch extremes \T{VLIW}.

\begin{bsp}[4-Tap-FIR auf $4\times4$-CGRA]
\label{sec:cgra}
$y[n]=h_0x[n]+h_1x[n{-}1]+h_2x[n{-}2]+h_3x[n{-}3]$,
Nearest-Neighbor-Mesh, 1-Takt-ADD/MUL.
\[m_k=h_k\,x[n{-}k],\quad a_0=m_0{+}m_1,\ a_1=m_2{+}m_3,\ y=a_0{+}a_1\]
\T{7 von 16 PE} benutzt $\;\;$ \T{1 Ergebnis pro Takt} (Pipeline)
$\;\;$ Latenz $=1_{\text{MUL}}+2_{\text{ADD}}=\T{3 Takte}$.
\end{bsp}

\SW{Xilinx AI Engines} Kachel aus Vektorprozessor (SIMD für parallele
Arithmetik) $+$ Skalarprozessor (Kontrollfluss, Adressen) $+$ lokalem
deterministischem Speicher $+$ DMA-Engines, verbunden über streaming-
und memory-mapped Interconnect.

\subsection{RP2350 PIO}
3 identische PIO-Blöcke, je \T{4 State Machines}. Realisiert in Software
Schnittstellen wie I2C, 3-pin I2S, SDIO, SPI/DSPI/QSPI, UART, DPI,
8080/6800-Parallelbus, VGA (Widerstands-DAC).
\begin{ztab}{@{}lX@{}}
Register je SM & $2\times32$\,Bit Shift (\T{OSR}, \T{ISR}),
  $2\times32$\,Bit Scratch (\T{X}, \T{Y}) \\
FIFO & $4\times32$\,Bit je TX und RX \\
Instruktionsspeicher & 32 Worte, gemeinsam für die 4 SM \\
Takt & fraktionaler Teiler (16 Integer-, 8 Nachkommabits) \\
Rest & flexibles GPIO-Mapping, DMA-Interface, IRQ-Flags \\
\end{ztab}

\begin{formel}[Nur 9 Befehle]
\R{JMP} \R{WAIT} \R{IN} \R{OUT} \R{PUSH} \R{PULL} \R{MOV} \R{IRQ}
\R{SET}. \R{NOP} ist eine Pseudo-Instruktion ($=$ \R{mov y,y}).\\
Format: \R{<instr> (side <wert>) ([<delay>])}, Delay 0--31 Takte
(weniger, wenn \R{.side\_set} benutzt wird).\\
JMP-Bedingungen: immer, \R{!X}, \R{X--}, \R{!Y}, \R{Y--}, \R{X!=Y},
\R{PIN}, \R{!OSRE}.
\end{formel}

\SW{Datenfluss} \R{PULL} holt ein Wort aus der TX-FIFO in den \T{OSR},
\R{OUT} schiebt 1--32\,Bit daraus zum Ziel; \R{IN} schiebt in den
\T{ISR}, \R{PUSH} schreibt ihn in die RX-FIFO.
\T{Autopull/Autopush} füllt bzw.\ leert automatisch
\ra die \R{pull}/\R{push}-Instruktion entfällt.

\SW{Zwei Tricks, die Takte sparen}
\begin{description}
\item[Side-Set] schaltet \emph{parallel} zum Befehl Pins um
  \ra Daten \emph{und} Takt ohne Zusatzbefehl (SPI).
\item[Program Wrapping] \R{.wrap\_target}/\R{.wrap} ersetzt den \R{jmp}
  durch einen Rücksprung, der \emph{null} Takte kostet.
\end{description}

\begin{lstlisting}
.program squarewave_wrap
  set pindirs, 1     ; Pin auf Ausgang
.wrap_target         ; 0-Takt-Ruecksprungziel
  set pins, 1 [1]    ; high, 1 Takt warten
  set pins, 0 [1]    ; low,  1 Takt warten
.wrap                ; 0-Takt-Sprung statt jmp
\end{lstlisting}

\SW{Direktiven} \R{.program}, \R{.define}, \R{.side\_set <n> (opt)},
\R{.in}/\R{.out <n> (left|right) (auto) (<schwelle>)},
\R{.fifo}, \R{.clock\_div}, \R{.mov\_status},
\R{.wrap\_target}/\R{.wrap}, \R{.origin}, \R{.pio\_version}
(1 $=$ RP2350, 0 $=$ RP2040), \R{.word}.

  % --- Werkzeuge ---
  % !TEX root = ../main.tex
% ---------------------------------------------------------------------
%  Docker -- Software-Virtualisierung mit Containern
% ---------------------------------------------------------------------
\section{Docker}
\label{sec:docker}

\SW{Wozu} Plattform zur \T{Software-Virtualisierung}: Software wird in
standardisierte, wiederverwendbare Einheiten verpackt, die \emph{alles}
enthalten, was zum Ausführen nötig ist. Diese Einheiten laufen auf nahezu
jedem Rechner direkt \ra portabel, skalierbar, unkompliziert.
Das Problem, das sie löst: „Bei mir läuft's!“ -- Abhängigkeiten,
Versionskonflikte, Installationsprozeduren.

\SW{Abgrenzung} Docker virtualisiert \emph{nur den Linux-Kernel}, nicht
die ganze Maschine. Deshalb gilt immer \T{Container-OS $=$ Host-OS} --
anders als bei einer virtuellen Maschine. Auf Windows/macOS braucht Docker
deshalb eine VM (WSL). Die Programme laufen direkt auf dem Host-Kernel;
in \R{ps ax} auf dem Host sieht man den Container-Prozess.

\subsection{Image vs.\ Container}
\begin{ztab}{@{}lXX@{}}
 & Image & Container \\ \midrule
Was & Bauplan, alles Nötige zum Ausführen & ein \emph{ausgeführtes} Image \\
Dateisystem & \T{read-only}, unveränderlich & \T{read-write} Overlay über dem Image \\
Anzahl & eines & beliebig viele aus einem Image \\
Platz & volle Grösse & nur die \emph{Änderungen} gegenüber dem Image \\
Lebensdauer & bis \R{rmi} & bis \R{rm}, überlebt das Stoppen \\
\end{ztab}

\begin{formel}[Der Satz, der alles trägt]
Das Ausführen eines \emph{Images} erzeugt einen \emph{Container} --
nie umgekehrt. Änderungen zur Laufzeit (Logfiles, DB-Einträge,
nachinstallierte Pakete) landen im Overlay des Containers, das Image
bleibt unberührt.
\end{formel}

\subsection{Host und Kapselung}
\SW{Host} der Rechner mit dem OS, auf dem die Container laufen (hier
Raspberry Pi OS). Er stellt nur die grundlegendsten OS-Funktionen bereit:
nacktes Dateisystem, Prozesse, Sockets. \emph{Nicht} Software, Daten,
User/Gruppen oder IP-Ports -- das muss der Container selbst mitbringen.

\SW{Kapselung} Aus dem Container heraus normalerweise kein Zugriff auf das
Dateisystem des Hosts oder anderer Container; gilt auch für User und Ports.
Der Host dagegen kann Container und Images ändern und löschen.
Die zwei Löcher in dieser Wand sind \R{docker cp} und \T{Volumes}.

\subsection{Lebenszyklus}
Vier Zustände: \T{created} $\to$ \T{running} $\to$ \T{stopped} $\to$
\T{deleted}, dazu das lokale Image-Repository als fünfter Knoten:
\begin{lstlisting}
Docker Hub --pull--> lokales Image-Repo
  lokales Repo --run-----> running
  lokales Repo --create--> created --start--> running
  running  --stop--> stopped --start/restart--> running
  stopped/created --rm--> deleted
  Container --commit--> lokales Image-Repo
\end{lstlisting}
\begin{ztab}{@{}lX@{}}
\R{docker run <img>} & Image \ra Container erzeugen \emph{und} starten \\
\R{docker create <img>} & nur erzeugen (Zustand \emph{created}) \\
\R{docker start <cnt>} & starten, \emph{detached} (Hintergrund) \\
\R{docker start -i <cnt>} & starten, \emph{attached} (Vordergrund) \\
\R{docker stop <cnt>} & running \ra stopped \\
\R{docker restart <cnt>} & stop, dann start \\
\R{docker rm <cnt>} & Container löschen \\
\R{docker rmi <img>} & Image löschen \\
\end{ztab}

\begin{falle}[Drei Stolperstellen]
1.~\T{Jedes \R{run} erzeugt einen \emph{neuen} Container} -- nicht
\R{run} zum Wiederstarten nehmen, sondern \R{start}.\\
2.~Ein Image lässt sich \textbf{nicht} löschen, solange noch ein Container
davon existiert.\\
3.~Wer \R{-it} beim \R{create} vergisst, kann den Container später nicht
interaktiv nutzen -- die Option muss schon beim \emph{Erzeugen} dabei sein.
\end{falle}

\subsection{Die wichtigsten Kommandos}
\begin{ztab}{@{}lX@{}}
\R{docker image ls} & lokale Images (kurz: \R{docker images}) \\
\R{docker ps} & laufende Container (lang: \R{container ls}) \\
\R{docker ps -a} & auch beendete \\
\R{docker ps -a --size} & mit Platzbedarf \\
\R{docker run -it <img>} & interaktiv (siehe unten) \\
\R{docker run --rm <img>} & Container nach Ende sofort löschen \\
\R{docker run --name <n> <img>} & eigener Name statt Zufallsname \\
\R{docker rename <alt> <neu>} & nachträglich umbenennen \\
\R{docker attach <cnt>} & Terminal an Hintergrund-Container hängen \\
\R{docker container prune} & \emph{alle} gestoppten Container löschen \\
\R{docker login -u <user>} & an Docker Hub anmelden \\
\R{docker image push <tag>} & Image auf Docker Hub laden \\
\end{ztab}
Docker vergibt jedem Container eine zufällige \T{ID} und einen zufälligen
\T{Namen}; Kommandos akzeptieren beides. Fehlt ein Image lokal, lädt
\R{run} es automatisch von Docker Hub.

\begin{formel}[Was \R{-i} und \R{-t} bedeuten]
\R{-i} verbindet den Container mit der Standardeingabe, \R{-t} legt einen
Pseudo-Terminalemulator davor. Zusammen \R{-it} ergibt eine
Terminal-Sitzung wie bei \R{ssh}.\\
\T{Ohne \R{-it} beendet sich ein Shell-Container sofort}: die \R{bash}
findet keine Standardeingabe und bricht ab -- der häufigste Anfängerfehler.
\end{formel}

\subsection{Beenden und Trennen}
\begin{ztab}{@{}lX@{}}
\R{exit} & beendet die Shell im Container \ra Container \emph{stopped} \\
\R{<Ctrl-d>} & dasselbe (beendet die Shell) \\
\R{<Ctrl-c>} & \emph{nicht-interaktiv}: beendet \R{docker run} und damit
  den Container; \emph{interaktiv}: beendet nur das gerade laufende
  Programm \emph{in} der Shell, nicht die Shell \\
\R{<Ctrl-p> <Ctrl-q>} & \T{detach}: nur das Terminal trennen, Container
  läuft weiter. Nur bei Containern mit \R{-it}. Zurück mit \R{attach}. \\
\end{ztab}
In allen Fällen \emph{existiert} der Container weiter (Status
\emph{Exited}), bis man ihn mit \R{rm} löscht.

\subsection{Schlanke Images}
Images sind bewusst minimal -- im Ubuntu-Image fehlen \R{ping}, \R{nano},
\R{wget}, \R{curl}; vorhanden sind \R{cp}, \R{mv}, \R{mkdir}, \R{cat},
\R{tar}, \R{apt}, \R{ps}, \R{top}.
\SW{Nachinstallieren} erst die Paketlisten aktualisieren, dann installieren:
\begin{lstlisting}
apt -y update && apt -y install nano
\end{lstlisting}
\R{-y} beantwortet alle Rückfragen mit Ja (nötig, weil beim Bauen von
Images niemand tippt). Die Verkettung mit \R{\&\&} ist üblich -- warum,
zeigt sich beim Dockerfile.

\SW{Auch die Prozesse sind minimal} Auf einem normalen Linux laufen viele
Dienste (\R{sshd}, \R{dockerd}, \R{dhcpd}, \dots{} -- in Unix
\emph{daemons}, Endung \R{d}). Im Ubuntu-Container läuft \T{genau ein
Prozess: \R{bash}} -- der Container soll nur eine Software ausführen, den
Rest stellt der Host bereit. Zählen mit \R{ps -aux | wc -l}.

\begin{bsp}[Platzbedarf -- typische Prüfungsfrage]
\R{docker ps -a --size} zeigt \R{0B (virtual 192MB)}: der Container selbst
belegt nichts, solange er nichts geändert hat; die 192\,MB sind das Image.
Installiert man darin 20\,MB Software, wächst der \emph{Container} um
20\,MB, das Image bleibt gleich gross.\\
\ra 4 Container aus einem 50-MB-Image ohne Änderungen brauchen zusammen
\T{50\,MB}, nicht 200. Und: Dateien \emph{im Container} zu löschen gibt
keinen Platz frei, wenn sie aus dem Image stammen.
\end{bsp}

\subsection{Hintergrund und zusätzliche Prozesse}
\begin{ztab}{@{}lX@{}}
\R{docker run -d <img>} & \T{detached} starten; läuft, bis der gestartete
  Prozess sich beendet. Für Serverdienste ohne Shell-Interaktion. \\
\R{docker run -itd <img>} & interaktiv \emph{und} detached; kehrt sofort
  zurück, später \R{attach} \\
\R{docker logs <cnt>} & Ausgabe des Containers ansehen \\
\R{docker logs -f <cnt>} & \emph{folgt} der Ausgabe wie \R{tail -f};
  Abbruch mit \R{<Ctrl-c>} \\
\R{docker exec <cnt> <cmd>} & \T{zusätzlichen} Prozess im laufenden
  Container starten \\
\R{docker exec -it <cnt> bash} & zweite Shell parallel -- das
  Debug-Werkzeug \\
\R{docker exec -d <cnt> <cmd>} & dieser Prozess im Hintergrund \\
\R{docker top <cnt> ax} & Prozesse \emph{im} Container auflisten \\
\end{ztab}

\begin{formel}[\R{run} gegen \R{exec} -- der Unterschied]
\R{run} erzeugt einen \T{neuen Container} aus einem Image.
\R{exec} startet einen \T{weiteren Prozess} in einem \emph{bereits
laufenden} Container. \R{exec} auf einen gestoppten Container
funktioniert nicht.
\end{formel}

\subsection{Dateien kopieren: \R{docker cp}}
\begin{lstlisting}
docker cp <cnt>:<quelle> <ziel>    # Container -> Host
docker cp <quelle> <cnt>:<ziel>    # Host -> Container
\end{lstlisting}
Funktioniert auch bei einem Container im Zustand \emph{created}, also
bevor er je lief. Typische Abfolge: \R{docker create --name myubuntu -it
ubuntu} $\to$ \R{docker cp hello.txt myubuntu:/tmp/} $\to$
\R{docker start myubuntu} $\to$ \R{docker attach myubuntu}.

\subsection{Das geschichtete Container-Dateisystem}
Das virtuelle Dateisystem eines Containers besteht aus übereinander
gelegten Schichten:
\begin{ztab}{@{}lX@{}}
oben & \T{Container-Layer}, read-write -- die Änderungen \\
Mitte & weitere \emph{Image}-Schichten: nachinstallierte Software,
  Webserver, DB, dotnet \dots{} \\
unten & \T{Base-Image}, meist ein minimales Linux (Ubuntu, Debian) \\
\end{ztab}
Alle Schichten ausser der obersten sind \T{read-only}.
\T{Jeder Container bekommt seinen eigenen Container-Layer} -- auch mehrere
aus demselben Image. Der Layer wird \emph{mit} dem Container gelöscht.
\merkzeile Praktisch für reinen Tisch (temporäre Dateien verschwinden),
unpraktisch für Software, die Daten dauerhaft speichern muss.

\SW{Wo das physisch liegt} \R{/var/lib/docker/overlay2/\dots} und
\R{/var/lib/containerd/\dots}. Suchen mit
\R{sudo find /var/lib/ -name <datei>}.

\begin{falle}[Zwei Beobachtungen dazu]
1.~Bei einem \emph{laufenden} Container findet \R{find} die Datei
\T{zweimal} (docker und containerd), bei einem gestoppten nur einmal.\\
2.~Diese Verzeichnisse werden von Docker verwaltet und dürfen
\textbf{nie} direkt verändert werden.
\end{falle}

\subsection{Volumes}
Ein Volume bildet einen Pfad des Host-Dateisystems in das
Container-Dateisystem ab. Änderungen sind auf beiden Seiten sofort
sichtbar, und die Dateien \T{überleben das Löschen des Containers} --
alles andere im Container-Layer nicht.

\begin{ztab}{@{}lX@{}}
\T{Host Volume} & \R{-v /pfad/im/host:/pfad/im/cnt} -- Nutzer bestimmt
  beide Seiten \\
\T{Anonymes V.} & \R{-v /pfad/im/cnt} -- Docker legt das Host-Verzeichnis
  unter \R{/var/lib/docker/volumes/<hash>/\_data/} an; kaum gebraucht \\
\T{Named Volume} & \R{-v name:/pfad/im/cnt} -- wie anonym, aber mit Namen \\
\end{ztab}

\begin{bsp}[Wofür Volumes wirklich da sind]
Versions-Update eines Datenbankservers: Die containerisierte DB legt ihre
Daten im Volume ab. Beim Update wird der alte Container gestoppt und
gelöscht -- die Daten bleiben. Der neue Container aus dem aktuellen Image
greift auf dieselben Daten zu.
\end{bsp}

\SW{Administration} \R{docker volume ls} $\cdot$
\R{docker volume rm <vol>} $\cdot$ \R{docker volume prune}.
Volumes überleben ihren Container, müssen also von Hand aufgeräumt werden.

\begin{falle}[Warum \R{volume prune} trügerisch ist]
Es löscht \emph{alle} Volumes, die gerade keinem Container zugeordnet
sind. Genau die wertvollen Daten eines Volumes, dessen Container du
versehentlich schon gelöscht hast, sind damit weg -- und es gibt
\textbf{kein} Docker-Kommando, um vorher den Inhalt eines Volumes
anzusehen.
\end{falle}

\subsection{Image aus Container: \R{docker commit}}
\begin{lstlisting}
docker commit <container> <image>
\end{lstlisting}
Friert die Änderungen eines laufenden oder gestoppten Containers als neue
Image-Schicht ein. Damit lässt sich manuell installierte Software in ein
Image übernehmen: Ubuntu starten $\to$ \R{apt install nano} $\to$
\R{commit} $\to$ das neue Image enthält nano.
\merkzeile Das ist die Antwort auf „Kann aus einem Container ein Image
erzeugt werden?“ -- ja, mit \R{commit}. Der saubere Weg bleibt aber das
Dockerfile.

\subsection{Repositories und Docker Hub}
\T{Docker Hub} ist das Standard-Repository: \R{docker run} und
\R{docker pull} holen Images von dort. Öffentliche Repositories sind für
jeden zugänglich, private nur für den Besitzer (bei Docker Hub ab einer
gewissen Anzahl kostenpflichtig) -- oder on premise mit eigener
Repository-Software (CNCF Distribution).

\begin{ztab}{@{}lX@{}}
\R{docker search <begriff>} & Images suchen; Spalten \R{STARS},
  \R{OFFICIAL}, \R{AUTOMATED}. Offizielle Images sind vertrauenswürdiger. \\
\R{docker pull <img>} & herunterladen, ohne zu starten \\
\end{ztab}

\begin{formel}[Aufbau eines Image-Namens]
\begin{center}\ttfamily
<repo>/<source>/<image>:<tag>
\end{center}
\begin{tabularx}{\linewidth}{@{}lX@{}}
\R{<image>} & der Name, z.\,B. \R{ubuntu}, \R{hello-world} \\
\R{<source>} & Nutzer oder Organisation auf Docker Hub;
  Standard \R{library} (die offiziellen Images) \\
\R{<tag>} & die Version; Standard \R{latest} \\
\R{<repo>} & die Registry; Standard \R{index.docker.io} \\
\end{tabularx}\\[2pt]
Also ist \R{docker run hello-world} identisch mit\\
\R{docker run index.docker.io/library/hello-world:latest}.\\
Weitere Beispiele: \R{ubuntu:jammy} (22.04),
\R{kaohslu/epp-demo}, \R{mcr.microsoft.com/dotnet/runtime:5.0}.
\end{formel}

\begin{falle}[\R{run} aktualisiert nicht]
\R{docker run <img>} nimmt die \T{neueste lokal vorhandene} Version --
auch wenn im Repository längst eine neuere liegt. Nur \R{docker pull}
holt die aktuelle.\\
Und beim \R{pull}: Das alte Image wird \emph{nicht} überschrieben,
sondern verliert nur seinen Tag und erscheint mit \R{<none>}
(\emph{dangling image}) -- Container könnten es noch brauchen.
\end{falle}

\subsection{Der Weg eines eigenen Images}
\begin{lstlisting}
docker build -t <user>/<name> .   # Image bauen und taggen
docker login -u <user>            # an Docker Hub anmelden
docker image push <user>/<name>   # hochladen
docker run <user>/<name>          # auf beliebigem Host starten
\end{lstlisting}
Drei Zeilen zum Veröffentlichen, eine zum Benutzen -- ohne dass auf dem
Zielrechner irgendeine Laufzeitumgebung installiert sein muss.

  % --- Referenz ---
  % !TEX root = ../main.tex
% ---------------------------------------------------------------------
%  Referenz: Formeln, Rezepte, durchgerechnete Beispiele
% ---------------------------------------------------------------------
\section{Formeln \& durchgerechnete Beispiele}
\label{sec:referenz}

\subsection{Alle Formeln auf einen Blick}
\begin{ztab}{@{}lX@{}}
Adresszerlegung & $\#\text{Sets}=\frac{\text{Cache}}{\text{Line}\cdot\text{Assoz.}}$,
  Offset $=\log_2(\text{Line})$, Index $=\log_2(\#\text{Sets})$ \\
AMAT & $t_{\text{Hit}}+\text{MissRate}\cdot\text{MissPenalty}$ \\
Amdahl & $S=\dfrac{1}{(1-p)+p/s}$ \quad ($s$: Lanes \emph{oder} Kerne) \\
Bandbreite & Takt $\times$ Busbreite in Byte \\
Trashing-Abstand & $k\cdot\frac{\text{Cachegrösse}}{\text{Assoziativität}}$ \\
Page Table 1-stufig & $2^{\text{VPN-Bits}}\times$ Eintragsgrösse \\
PL330-Burst & $n_{\text{Beats}}\times m_{\text{Bit}}/8$ Byte je \R{DMALD} \\
CGRA-Latenz & Summe der Stufen entlang des Datenpfads \\
\end{ztab}

\subsection{Durchgerechnete Beispiele}
\begin{bsp}[Sieben Rechnungen, wie sie in der Prüfung drankommen]
\T{(a) Adresszerlegung.} 32\,kB Cache, 32\,B Lines, 8-way, 32-Bit-Adresse.
$\#\text{Sets}=\frac{32768}{32\cdot8}=128$ \ra Offset $=\log_232=5$\,Bit
(A4:A0), Index $=\log_2128=7$\,Bit (A11:A5), Tag $=32-7-5=\T{20\,Bit}$.

\T{(b) Trashing konstruieren.} Direct mapped, 8\,kB, 64\,B Lines
\ra 128 Lines, Index $=$ A12:A6. Kollision, wenn der Abstand ein
Vielfaches der Cachegrösse ist: \R{0x0000} und \R{0x2000} -- gleicher
Index, verschiedene Tags.

\T{(c) AMAT zweistufig.} $t_{L1}=4$, MissRate$_{L1}=10\,\%$, $t_{L2}=12$,
MissRate$_{L2}=25\,\%$, RAM $=120$:
$4+0.1\,(12+0.25\cdot120)=\T{8.2 Cycles}$.

\T{(d) Amdahl mit Datentypwahl.} 90\,\% vektorisierbar.
\R{float} (4 Lanes): $1/(0.1+0.9/4)=\T{3.08}\times$;
\R{uint8} (16 Lanes): $1/(0.1+0.9/16)=\T{6.4}\times$.
\ra Die Datentypwahl verdoppelt den Speedup, der Code bleibt gleich.

\T{(e) PL330-CCR.} \R{SB8 SS32 DB4 DS64} mit \R{DMALP 32}: Quelle
$8\cdot32/8=32$\,B, Ziel $4\cdot64/8=32$\,B -- passt. Gesamt
$32\cdot32=\T{1024\,B}$. Schleifenkörper 4 Instruktionen \ra passt in
eine I-Cache-Line.

\T{(f) \R{TI} für UART-TX} (mem\,\ra\,Peripherie, 32\,Bit):
\R{SRC\_INC}$=1$, \R{DEST\_INC}$=\mathbf{0}$, \R{DEST\_DREQ}$=1$,
\R{SRC\_DREQ}$=0$, \R{PERMAP}$=$ UART-Nr., \R{WIDTH}$=32$, \R{INTEN}$=1$.

\T{(g) Zynq-DMA-Bandbreite.} CPU\_6x $=$ 800\,MHz, 1:3 \ra CPU\_2x $=$
266\,MHz, Interconnect 64\,Bit $=$ 8\,B: $266\cdot8=\T{2.13\,GB/s}$
(Slave-Interconnect 32\,Bit: $\sim$1.07\,GB/s).
\end{bsp}

\subsection{Die vier durchgehenden Muster}
\begin{description}
\item[Overhead amortisieren] Cache Line, AXI-Burst, SIMD-Lane, Loop
  Unrolling -- viermal dasselbe Argument: eine feste Startgebühr auf
  möglichst viele Nutzdaten verteilen.
\item[Kleiner schneller Speicher vor grossem langsamem] L1/RAM,
  TLB/Page Table, PL330-I-Cache -- immer dieselbe Struktur (Tag, Valid,
  Miss, Replacement), nur die Zeitskalen verschieben sich um Faktor 1000.
\item[Prädikation statt Verzweigung] NEON-Masken (räumlich, über Lanes),
  PL330-\R{B}/\R{S}-Suffixe (zeitlich, über Instruktionen).
\item[Die Kette] P1 baut die Hierarchie $\;\to\;$ P2 macht die CPU
  schneller und verschärft damit das Speicherproblem $\;\to\;$ P3 nimmt
  der CPU das Datenschaufeln ab $\;\to\;$ P4 vereinfacht den Kern
  $\;\to\;$ P5 vervielfacht ihn und zahlt mit Kohärenz $\;\to\;$ P6
  verlagert die Arbeit ganz aus der CPU heraus.
\end{description}

  % !TEX root = ../main.tex
% ---------------------------------------------------------------------
%  Stichwortregister -- vom Begriff zur Abschnittsnummer
%  ERZEUGT von mkreg.py. Neue Eintraege dort ergaenzen und neu laufen
%  lassen, oder hier von Hand alphabetisch einsortieren.
% ---------------------------------------------------------------------
\section{Stichwortregister}

\scriptsize\renewcommand{\arraystretch}{0.9}
\begin{tabularx}{\linewidth}{@{}Xr@{\hskip4pt}Xr@{\hskip4pt}Xr@{}}
\R{-it} & \ref{sec:docker} & Doorbell & \ref{sec:multicore} & PIO & \ref{sec:reconfig} \\
Adresszerlegung & \ref{rez:zerlegung} & DREQ / TREQ & \ref{sec:dma} & PL330 & \ref{sec:dma} \\
AMAT & \ref{form:amat} & Excl. Monitor & \ref{sec:multicore} & Prädikation & \ref{sec:neon} \\
Amdahl & \ref{form:amdahl} & Flags (fehlen) & \ref{sec:rvflags} & Prio-Inversion & \ref{sec:kohaerenz} \\
apt (Container) & \ref{sec:docker} & \R{FPEXC} & \ref{sec:neon} & \R{prune} & \ref{sec:docker} \\
Assoziativität & \ref{sec:cache} & FPGA / LUT & \ref{sec:reconfig} & Race Condition & \ref{sec:kohaerenz} \\
Attach / Detach & \ref{sec:docker} & FreeRTOS-API & \ref{sec:multicore} & Reentranz & \ref{sec:multicore} \\
Autopull/-push & \ref{sec:reconfig} & Harvard & \ref{sec:cache} & Replacement & \ref{sec:cache} \\
AXI-Burst & \ref{sec:dma} & Image (Docker) & \ref{sec:docker} & Ringpuffer (HW) & \ref{sec:dma} \\
Bandbreite & \ref{form:bw} & Image-Name / Tag & \ref{sec:docker} & RISC-Prinzip. & \ref{sec:riscv} \\
Base-Image & \ref{sec:docker} & Index & \ref{rez:zerlegung} & Routing-Delay & \ref{sec:reconfig} \\
\R{BR}/\R{BG} & \ref{sec:dma} & Intrinsics & \ref{sec:neon} & RV32C & \ref{sec:riscv} \\
\R{BSWAP} & \ref{sec:dma} & Kohär. Cache & \ref{falle:dma} & Sättigung (Q) & \ref{rez:neondecode} \\
Cache Line & \ref{sec:cache} & Kohär. Multicore & \ref{sec:kohaerenz} & Scatter-Gath. & \ref{falle:dma} \\
CAM & \ref{sec:cache} & Lanes & \ref{sec:neon} & Semaphore & \ref{sec:kohaerenz} \\
\R{CCR} & \ref{rez:ccr} & \R{latest} & \ref{sec:docker} & Shape L/W/N & \ref{rez:neondecode} \\
CGRA & \ref{sec:cgra} & Lebenszyklus & \ref{sec:docker} & Side-Set & \ref{sec:reconfig} \\
\R{CHAIN\_TO} & \ref{sec:dma} & Locality & \ref{sec:cache} & SIO & \ref{sec:multicore} \\
Container & \ref{sec:docker} & Loop Peeling & \ref{sec:dma} & \R{SLTU} & \ref{sec:rvflags} \\
Container-Layer & \ref{sec:docker} & Loop Unrolling & \ref{sec:neon} & SMP / AMP & \ref{sec:multicore} \\
Control Block & \ref{rez:ti} & LRU & \ref{sec:cache} & Spinlock & \ref{sec:multicore} \\
Copy-on-write & \ref{sec:mmu} & MAC / \R{VMLA} & \ref{sec:neon} & \R{SRC\_INC} & \ref{rez:ti} \\
\R{CPACR} & \ref{sec:neon} & Mailbox & \ref{sec:kohaerenz} & Strip Mining & \ref{sec:rvv} \\
CRC im DMA & \ref{sec:dma} & Maske (\R{VBSL}) & \ref{sec:neon} & Swapping & \ref{sec:mmu} \\
Cycle Stealing & \ref{sec:dma} & Matrixmult. & \ref{sec:neon} & Tag & \ref{rez:zerlegung} \\
dangling image & \ref{sec:docker} & MFIFO & \ref{sec:dma} & Tail-Handling & \ref{sec:simdgrenzen} \\
detached (-d) & \ref{sec:docker} & MMU & \ref{sec:mmu} & \R{TI}-Register & \ref{rez:ti} \\
Dirty Bit & \ref{sec:cache} & MPU & \ref{sec:mmu} & TLB & \ref{sec:mmu} \\
\R{DMAWFP} & \ref{sec:dma} & Mutex & \ref{sec:kohaerenz} & Trashing & \ref{bsp:trash} \\
\R{docker commit} & \ref{sec:docker} & Named Volume & \ref{sec:docker} & Valid Bit & \ref{sec:cache} \\
\R{docker cp} & \ref{sec:docker} & \R{NEXTCONBK} & \ref{rez:ti} & VALID/READY & \ref{sec:dma} \\
\R{docker exec} & \ref{sec:docker} & NMRU & \ref{sec:cache} & Vektorlänge VL & \ref{sec:rvv} \\
Docker Hub & \ref{sec:docker} & Offset & \ref{rez:zerlegung} & \R{VLD3} & \ref{sec:neon} \\
\R{docker logs} & \ref{sec:docker} & Overlay-FS & \ref{sec:docker} & \R{vm\_stat} & \ref{sec:mmu} \\
\R{docker ps} & \ref{sec:docker} & Page Fault & \ref{sec:mmu} & Volume & \ref{sec:docker} \\
\R{docker pull} & \ref{sec:docker} & Page Table & \ref{sec:mmu} & \R{volume prune} & \ref{sec:docker} \\
\R{docker run} & \ref{sec:docker} & Pairwise & \ref{sec:neon} & \R{vsetvl} & \ref{sec:rvv} \\
\R{docker search} & \ref{sec:docker} & Panic-Signal & \ref{sec:dma} & Wrapping (PIO) & \ref{sec:reconfig} \\
\R{docker top} & \ref{sec:docker} & \R{PERMAP} & \ref{rez:ti} & Write-Back & \ref{sec:cache} \\
\end{tabularx}
\normalsize

\end{multicols*}

\end{document}
